\documentclass[11pt,oneside]{report}
\usepackage[margin=1in]{geometry}
\usepackage{amsmath,amssymb,amsthm}
\usepackage{booktabs}
\usepackage{array}
\usepackage[hidelinks]{hyperref}
% let \path break at underscores and dots too, not just slashes
\expandafter\def\expandafter\UrlBreaks\expandafter{\UrlBreaks\do\_\do\.}

\theoremstyle{plain}
\newtheorem{theorem}{Theorem}[chapter]
\newtheorem{proposition}[theorem]{Proposition}
\newtheorem{lemma}[theorem]{Lemma}
\newtheorem{corollary}[theorem]{Corollary}
\newtheorem{barrier}[theorem]{Barrier}
\newtheorem{openproblem}[theorem]{Open Problem}
\theoremstyle{definition}
\newtheorem{definition}[theorem]{Definition}
\theoremstyle{remark}
\newtheorem{remark}[theorem]{Remark}

\newcommand{\N}{N}
\newcommand{\g}{\mathfrak{g}}
\newcommand{\Ee}{E}

% Each chapter of this thesis is written to stand alone as a paper;
% the chapterabstract environment carries the per-chapter abstract.
\newenvironment{chapterabstract}%
  {\begin{quotation}\small\itshape\noindent}%
  {\end{quotation}\vspace{0.5em}}

\begin{document}

% ---------------------------------------------------------------- title page
\begin{titlepage}
\centering
\vspace*{4em}
{\Huge\bfseries Cute Goldbach Gaps\par}
\vspace{1.5em}
{\Large Extremal Least Goldbach Summands:\\ Constructions, Costs, and Closure\par}
\vspace{4em}
{\large A thesis on the engineering of Goldbach deserts\par}
\vspace{6em}
{\large Anonymous\par}
\vspace{1em}
{\large July 2026\par}
\vfill
\begin{minipage}{0.8\textwidth}\small\centering
Every integer, congruence system, primality certificate, verifier, and
search log referenced in this thesis is committed to the accompanying
artifact repository and can be re-checked from scratch on a laptop.
Nothing rests on a log we could not regenerate.
\end{minipage}
\vspace*{3em}
\end{titlepage}

% ---------------------------------------------------------------- abstract
\chapter*{Abstract}
\addcontentsline{toc}{chapter}{Abstract}

For an even integer $\N$, the \emph{least Goldbach summand} is
$\g(\N)=\min\{p\ \text{prime}:\N-p\ \text{prime}\}$ --- the quantity $p(n)$ of
the minimal Goldbach partition studied by Granville, van de Lune and te
Riele. Nature keeps it tiny: exhaustive verification gives
$\g(\N)\le9\,781$ for every even $\N\le4\cdot10^{18}$. This thesis is about
the opposite, extremal and constructive question: \emph{how large can
$\g(\N)$ be forced to be, and at what cost in the size of $\N$?}

Three record games organise the answer. The \emph{height game} maximises
$\g(\N)$ with no size limit; the \emph{threshold game} minimises $\N$
subject to $\g(\N)>Q$; the \emph{budget game} maximises $\g(\N)$ subject
to $\N<X$. After proving unconditionally that $\g$ is unbounded and
transferring the Ford--Green--Konyagin--Maynard--Tao long-gap machinery
into an explicit lower bound
$\g(\N)\gg\log\N\,\log_2\N\,\log_4\N/\log_3\N$ along even integers with
exhibited partitions, the thesis develops the partial-cover search
paradigm --- congruence covers of the prime offsets, Chinese-remainder
progressions, and probable-prime scanning --- into a measured, predictive
science with a validated cost model, and uses it to set every record it
then studies. The certified board: $\g(\N)=1\,157\,341$ at $2\,480$
digits for the height game, exceeding the longest prime gap ever
certified while using integers seven times shorter than that gap's own
endpoints; a $150$-digit $\N$ with $\g(\N)=104\,527$ for the threshold
game $T(100\,000)$, forty-nine digits below the first-generation record;
and $\g(\N)=119\,419$ below $10^{199}$ for the budget game. A companion
construction of $2\,692$ digits carries a Goldbach desert $126.5$ times
longer than its fully-certified local prime gap, separating the two
phenomena experimentally.

The second half of the thesis asks why the records stop where they do,
and answers with mathematics rather than fatigue. The search exponent is
governed by a density model validated to $\pm0.1$ nats across
$1.2\cdot10^9$ scanned candidates; the digit frontier at fixed threshold
is an explicit, measured curve; and beneath both lie provable barriers.
Compressing constructions below $\sim$100 digits reduces to modular
subset-sum with per-position choices at density $\approx1.026$ ---
the regime where lattice reduction and Wagner-style list merging both
fail --- and to Chinese-remainder list decoding beyond the Johnson
radius, which practical lattice reduction does not penetrate even on
solution-abundant instances (measured). A closure theorem for
compositeness-certificate schemes then shows, across a concrete grammar
of algebraic mechanisms swept by $13\,661$ hand-enumerated and $41$
LLM-evolved instances, that congruence covering is the unique
economical mechanism: every alternative family in the grammar is closed
by a proof of its own kind, and the game's whole exponent is an
integrality gap that survives exact optimization over half the cover at
once. What remains open is named precisely: a rounding-gap programme
importing the Erd\H{o}s--Rankin covering technology constructively, a
sub-modexp compositeness-detection question posed to complexity theory,
two coding-theoretic problems, and a two-sided construction that would
uncap the height game by making the partition's primality certificate
free. The thesis closes with the walls each game will hit under any
plausible compute, and the records' honest epitaph: upper bounds,
manufactured, falling toward a floor they will never touch.

% ---------------------------------------------------------------- preface
\chapter*{Preface: how to read this thesis}
\addcontentsline{toc}{chapter}{Preface}

Each chapter is written to stand alone as a paper, with its own abstract,
and the chapters are ordered as the research was: construction first,
measurement second, impossibility last.

\begin{itemize}
\item \textbf{Chapter~\ref{ch:intro}} defines the least Goldbach summand,
  the empirical background, and the three record games, and states the
  final record board.
\item \textbf{Chapter~\ref{ch:foundations}} is the foundations paper:
  factorial and Wilson obstructions, primorial deserts, unconditional
  unboundedness, prime-offset congruence covers, and the transfer of the
  large-prime-gap lower bounds. It subsumes, updates, and supersedes the
  July 2026 manuscript \emph{Cute Goldbach Gaps}, whose first-generation
  records ($\g=109\,621$ at $237$ digits, $\g=105\,667$ at $199$ digits)
  are the baseline everything later improves.
\item \textbf{Chapter~\ref{ch:costmodel}} builds the quantitative theory
  of partial-cover search: the density model, its five-scale validation,
  the digit--compute exchange rate, the frontier curves, and the
  statistical discipline that kept negative results honest.
\item \textbf{Chapter~\ref{ch:records}} is the records paper: two
  independent pipelines, the search engine and its measured
  optimizations, the certification standard, and the certified ladder
  from $199$ digits down to $150$.
\item \textbf{Chapter~\ref{ch:gaps}} compares Goldbach deserts with
  ordinary prime gaps and reports the height-game records, including the
  desert that is $126.5$ times longer than its own certified prime gap.
\item \textbf{Chapter~\ref{ch:limits}} proves the compression barriers:
  why the threshold record cannot be pushed below $\sim$100 digits by
  any known technique, with the failed attacks measured rather than
  presumed.
\item \textbf{Chapter~\ref{ch:closure}} states and proves (at draft
  rigor, on a concrete grammar) the Divisor-Paradigm Closure Theorem:
  within a language of planted compositeness certificates, nothing beats
  covering --- and reports the adversarial falsification engine that
  tried to break it, including the accounting hole it found in our own
  referee first.
\item \textbf{Chapter~\ref{ch:conclusions}} collects what is proved,
  what is measured, what is conjectured, and the open problems, ranked.
\item \textbf{Appendix~\ref{app:artifacts}} maps every claim to its
  machine-checkable artifact and gives the verification runbook.
\end{itemize}

Three writing conventions, kept throughout. First, notation is explained
in words at first use --- the intended reader is a mathematically literate
non-specialist, and no step should require outside references to follow.
Second, every record claim separates three obligations: \emph{(a)} no
small prime is a summand, \emph{(b)} a Goldbach partition exists,
\emph{(c)} the least summand equals the stated value; witnesses for (a)
are unconditional, and (b) always means an exhibited partition with a
certified prime complement, never an appeal to a conjecture. Third,
negative results --- campaigns that found nothing, attacks that failed,
one retracted claim --- are reported with the same care as the records,
because the thesis's second half is built out of them.

The research was carried out in 2026 across several cooperating
sessions --- cover optimization and certification on 4-core CPU
containers, a GPU search fleet on ephemeral cloud notebooks, and
delegated agent sessions for audits and sweeps --- whose merged artifact
history is preserved in the repository. Where a result was produced by
one pipeline and re-verified by the other, the text says so, because that
independence is load-bearing for trust.

\tableofcontents

\chapter{Introduction: three games and one exponent}\label{ch:intro}

Goldbach's conjecture asserts that every even integer $\N>2$ is a sum of two
primes; it has been verified for all $\N\le4\cdot10^{18}$~\cite{OeS}.
Heuristically one expects not just one representation but many. The
Hardy--Littlewood circle method---which estimates such counts by treating the
primes as a random set of the appropriate density and then correcting for
divisibility constraints---predicts~\cite{HL}
\begin{equation}\label{eq:HL}
r(\N)\;\sim\;2\,\Pi_2\,\frac{\N}{(\log \N)^2}\prod_{\substack{p\mid \N\\ p>2}}\frac{p-1}{p-2},
\qquad
\Pi_2=\prod_{p>2}\Bigl(1-\tfrac{1}{(p-1)^2}\Bigr),
\end{equation}
where $r(\N)=\#\{p\le\N:p\text{ and }\N-p\text{ both prime}\}$. Reading this
from left to right: $\N/(\log\N)^2$ is the naive guess, since a number near
$\N$ is prime with probability about $1/\log\N$ and there are $\N$ choices for
$p$; the twin-prime constant $\Pi_2\approx0.6601$ corrects for the fact that
primes are not truly random modulo small numbers; and the last factor inflates
the count when $\N$ is highly composite. The conclusion is that a typical large
even integer should have \emph{many} Goldbach partitions. So the interesting
question is not whether a representation exists, but how lopsided the most
lopsided one must be.

\begin{definition}
For an even integer $\N\ge4$ admitting a Goldbach representation, the
\emph{least Goldbach summand} is
\[
\g(\N)\;=\;\min\{\,p\ \text{prime}:\N-p\ \text{is prime}\,\}.
\]
\end{definition}

This is the quantity $p(n)$ of the \emph{minimal Goldbach partition} studied by
Granville, van de Lune and te Riele~\cite{GLvdtR}, who conjectured a ceiling
\begin{equation}\label{eq:GLR}
\g(\N)\;=\;O\bigl((\log \N)^2\log\log \N\bigr).
\end{equation}
Empirically $\g$ is tiny and grows very slowly; its record values are
catalogued in the OEIS~\cite{OEIS} (A020481, with records A025018/A025019).
Exhaustive search gives a slowly climbing ladder: $\g=1093$ at
$\N=60\,119\,912$; $\g(\N)\le5569$ for all even
$\N\le4\cdot10^{14}$~\cite{Richstein}; and the current exhaustive record
$\g=9781$, first attained at $\N=3\,325\,581\,707\,333\,960\,528$, with
$\g(\N)\le9781$ for \emph{all} even $\N\le4\cdot10^{18}$~\cite{OeS}. Strikingly,
these records occur at integers of at most $19$ digits.

We ask the opposite, extremal and constructive question:
\begin{quote}\itshape
How large can $\g(\N)$ be forced to be, and at what cost in the size of $\N$?
\end{quote}
An even integer $\N$ for which $\g(\N)$ is large is a \emph{Goldbach desert}:
a number all of whose small prime offsets are blocked. Constructing deserts is
the subject of this thesis, and the reason the subject is interesting is that
every construction turns out to be governed by a single number---an exponent
that measures how much luck the construction still needs after all its
cleverness is spent.

\section{The three games}

There are three natural record questions, and they reward different
constructions. The first is the unbounded \emph{height game}: exhibit as large
a value of $\g(\N)$ as possible, with no limit on the size of $\N$. An ordinary
prime gap is one way to play, but it is a strictly stronger construction: a gap
must block \emph{every} integer offset, while a desert needs to block only the
\emph{prime} offsets. Prime-gap constructions are therefore a strict subclass of
the constructions available here---a point we return to in
\S\ref{sec:gaps} with a concrete two-orders-of-magnitude example.

For comparisons in which size matters we use two finite games.

\begin{definition}[Threshold and budget games]\label{def:games}
For $Q\ge2$ and $X\ge4$, define
\[
T(Q)=\min\{\,\N \text{ even, Goldbach-representable}:\ \g(\N)>Q\,\},
\qquad
R(X)=\max\{\,\g(\N):\ \N<X\,\}.
\]
The \emph{threshold game} fixes $Q$ and tries to lower $T(Q)$; the
\emph{budget game} fixes $X$ and tries to raise $R(X)$. ``Fewer than $200$
decimal digits'' means the strict budget $\N<10^{199}$.
\end{definition}

\begin{remark}[A naming caution]
The budget game has been written both as $R(10^{199})$ (``fewer than $200$
digits'') and as $R(10^{200})$ (``at most $200$ digits'') in the literature
this thesis builds on, and the two are genuinely different: our $200$-digit
incumbent with $\g=112\,249$ satisfies $\N<10^{200}$ but exceeds $10^{199}$ by
$3.4\%$, which is precisely why a $199$-digit integer with the smaller value
$\g=119\,419$ is the record for the stricter budget. Throughout we state the
bound explicitly with each record.
\end{remark}

\section{Certified incumbents}

Table~\ref{tab:records} lists the programme's certified incumbents at the time
of writing. Every entry has been verified twice: once by the pipeline that
produced it and once by an independent verifier written against the record file
rather than against the search code, including a re-derivation of every
per-offset witness and an independent check of the elliptic-curve primality
certificate.

\begin{table}[h]
\centering\small
\caption{Certified incumbents. ``Witnesses'' counts prime offsets excluded by
an explicit congruence divisor plus those excluded by a strong base-$2$
compositeness witness; in every case the positive side is an exhibited prime
complement with an ECPP certificate.}\label{tab:records}
\begin{tabular}{lrrl}
\toprule
Game & Digits of $\N$ & $\g(\N)$ & Construction\\
\midrule
Height (unbounded) & $2\,480$ & $1\,157\,341$ & cover at $Q=1.16\cdot10^6$, $t=14\,544$\\
Height, gap-dominating & $2\,692$ & $1\,134\,871$ & cover at $Q=1\,113\,138$, $k=14$\\
Threshold $T(100\,000)$ & $150$ & $104\,527$ & cover at $Q=100\,003$, $k=6.21\cdot10^{10}$\\
Budget $R(10^{199})$ & $199$ & $119\,419$ & cover at $Q\approx1.19\cdot10^5$, $t=1.46\cdot10^8$\\
Budget $R(10^{200})$ & $200$ & $112\,249$ & cover at $Q=107\,720$, $k=56\,770$\\
\bottomrule
\end{tabular}
\end{table}

Two of these deserve comment. The $2\,692$-digit entry is included because it
settles the comparison with ordinary prime gaps in the strongest possible way:
its $\g$ exceeds the largest \emph{certified} prime gap known
($1\,113\,106$), and it does so with an integer roughly $6.9$ times shorter
than the one realising that gap. Deserts really are cheaper than gaps, by a
measurable factor.

The $150$-digit threshold record is the endpoint of a \emph{ratchet}: the same
covering machinery, re-tuned and re-run, produced a descending ladder
$195\to179\to150$ digits at $Q\approx10^5$, banking a certified record at each
rung. Chapter~\ref{ch:engineering} argues that this ratchet discipline, rather
than any single clever cover, is the strategy that actually produces records
under a fixed compute budget.

\section{The exponent that governs everything}

The technical spine of the thesis is a single quantity. A partial congruence
cover fixes $\N$ inside an arithmetic progression $\N=\N_0+kM$ and disposes of
most prime offsets for \emph{every} $k$; a residual set $U$ of offsets survives
and must be knocked out by luck, one candidate $k$ at a time. Under the
standard heuristic, the probability that a given candidate succeeds is
$e^{-E}$ with
\begin{equation}\label{eq:E}
E\;=\;\frac{|U|\cdot\mathrm{boost}}{\log\N},
\qquad
\mathrm{boost}\;=\;2\prod_{r\in\mathcal R}\frac{r}{r-1},
\end{equation}
so that $e^{E}$ candidates must be examined per expected hit. We call $E$ the
\emph{failure exponent}. It is the currency of the entire subject: every design
decision---which moduli to use, which residues, how many offsets to leave
residual, what digit budget to spend---is a trade in units of $E$, and every
piece of engineering (faster tests, better sieving, selective testing, more
hardware) buys a fixed number of nats of $E$ and no more.

Equation~\eqref{eq:E} is a heuristic, so we test it rather than assume it.
Across three independent covers and $1.2\times10^{9}$ scanned candidates, the
measured exponent $\Ehat$ agrees with the predicted $E$ to within $0.1$ nats
(Chapter~\ref{ch:engineering}). That agreement is what licenses the rest of the
thesis: once the model is trusted, questions about feasibility become
arithmetic.

\section{What is new here}

Chapters~\ref{ch:foundations} and~\ref{ch:engineering} are largely expository
and engineering: they set out the classical constructions, the covering
formulation, the calibrated cost model, and the search engine, and they report
the certified records and---equally important---the campaigns that failed.
The contributions begin in Chapter~\ref{ch:frontier}.

\begin{enumerate}
\item \textbf{The difficulty is an integrality gap} (\S\ref{sec:lp}). The
natural linear-programming relaxation of the covering problem is
\emph{vacuous}: fractionally, the available moduli cover $100.00\%$ of the
prime offsets at every target we tested. Consequently the whole failure
exponent is the price of \emph{rounding}: naive independent rounding leaves
$15.6\%$ of offsets uncovered, adaptive greedy leaves $7.56\%$, and the
fractional optimum leaves none. Every absolute percentage point recovered is
worth about $2.7$ nats, i.e.\ a factor of $15$ in compute.

\item \textbf{The rounding gap resists optimisation} (\S\ref{sec:door1}). We
attack it with exact large-neighbourhood search (certified optimal by LP
duality on $21$ of $22$ windows), with belief propagation and decimation, and
with a simultaneous exact solve of half the cover at once. The greedy solution
$|U|=867$ is never beaten. A random-restart probe shows why: the landscape is a
field of mutually distant, locally-exact-stable basins, and the adaptive greedy
construction lands about $100$ offsets deeper than a typical one.

\item \textbf{Sub-$100$-digit deserts reduce to a hard problem}
(Chapter~\ref{ch:barriers}). Oversized covers would make sub-$100$-digit
deserts easy if their Chinese-remainder representatives could be made small;
that requirement is exactly a modular subset-sum problem with per-position
choices at density $\approx1.026$, the regime where neither lattice reduction
nor birthday methods work. We also verify against the source that the relevant
CRT list-decoding radius is a factor $5.2$ short of what the construction would
need, and confirm experimentally that lattice methods do not penetrate the
Johnson-type radius even when solutions are abundant.

\item \textbf{A closure theorem} (Chapter~\ref{ch:closure}). Within a concrete
language of certificate schemas, no construction beats covering. Every family
closes for its own quantifiable reason---a $\sqrt{Q}$ value-set ceiling, a
purchase plan running $2.5\times$ over the design-entropy budget, a cluster
bound, lattice isolation, or an empty intersection at record scale---and the
theorem is conditional on two named hypotheses, one of which we pose as an open
problem.

\item \textbf{A methodological finding} (\S\ref{sec:referee}). We stress-tested
the closure theorem with an adversarial evolutionary search against a
machine-checkable scorer. Its first success was against the \emph{referee}, not
the theorem: it found a $50\times$ accounting error in our evaluator and
exploited it within ten candidates. The audit protocol caught it before any
claim was made, and repairing the referee strengthened the mathematics.
\end{enumerate}

\section{Notation and conventions}

Throughout, $\log_k$ is the $k$-fold iterated logarithm (so
$\log_2\N=\log\log\N$), and each extra $\log$ grows \emph{extremely} slowly;
$\theta(x)=\sum_{p\le x}\log p$ is Chebyshev's function;
$\gamma\approx0.5772$ is Euler's constant; $f\gg g$ means $f\ge cg$ for some
$c>0$ and all large arguments; $\pi(x)$ is the prime-counting function; and
``BPSW'' denotes the Baillie--PSW probable-prime test~\cite{PSW,BW}. All
entropies, costs and exponents are measured in \emph{nats} (natural logarithm
units), because the natural unit of a congruence modulo $r$ is $\log r$ and the
natural unit of a search is $\log(\text{candidates})$.

We are careful to separate three claims about any proposed record, and the
reader should hold us to the distinction throughout:
\begin{itemize}
\item[(a)] no prime below the stated value is a Goldbach summand of $\N$;
\item[(b)] $\N$ admits a Goldbach representation at all;
\item[(c)] the least summand \emph{equals} the stated value.
\end{itemize}
A congruence or a failed probable-prime test settles (a) unconditionally---a
prime can never fail such a test, so negative results are never in doubt. Claim
(b) needs an exhibited prime complement, and (c) needs both together with a
primality certificate for that complement. Where a statement is heuristic,
conditional, or merely measured, we say so: results proved at the level of
rigour of this thesis are marked as such, statements resting on standard but
unproved analytic heuristics are flagged, and purely empirical claims are
reported with the size of the experiment that supports them.

\chapter{Foundations: from factorials to covers, unconditionally}
\label{ch:foundations}

\begin{chapterabstract}
This chapter is the foundations paper, subsuming and updating the July
2026 manuscript \emph{Cute Goldbach Gaps}. After an elementary factorial
warm-up, we prove unconditionally that $\g$ is unbounded: a primorial
obstruction combined with Dirichlet's theorem produces, for every prime
$P$, an even $\N$ with an explicit partition and
$\g(\N)=\mathrm{nextprime}(P)$. We then formulate the construction as a
congruence-cover problem on the prime offsets, in complete and partial
forms, prove the realisation theorem that upgrades any cover to an
actual Goldbach number, and transfer the
Ford--Green--Konyagin--Maynard--Tao long-gap covering into an
unconditional lower bound on the maximal order of $\g$. The chapter
closes with the first-generation records this machinery produced --- a
$1020$-digit complete-cover baseline, a $237$-digit integer with
$\g=109\,621$, and a $199$-digit integer with $\g=105\,667$ --- which are
the incumbents every later chapter improves.
\end{chapterabstract}

\section{Factorial neighbourhoods and Wilson obstructions}\label{sec:factorial}

\begin{lemma}[Factorial composites]\label{lem:fact}
For $n\ge2$ and $2\le k\le n$, the number $n!+k$ is composite. For
$n\ge4$ and $2\le k\le n$, the number $n!-k$ is also composite.
\end{lemma}
\begin{proof}
In both cases $k\mid n!$, so $k\mid(n!\pm k)$. For the plus side
$n!+k>k$, so $n!+k$ is a proper multiple of $k$. For the minus side,
$n!-k=k\bigl(n!/k-1\bigr)$ with $n!/k-1\ge2$ exactly when $n!\ge3k$,
which holds for $n\ge4$; it fails at $n=3,\,k=3$, where $3!-3=3$ is
prime.
\end{proof}

\begin{lemma}[Wilson obstructions]\label{lem:wilson}
Let $n\ge2$. \emph{(A)} If $n+1$ is prime, then $(n+1)\mid n!+1$.
\emph{(B)} If $n+2$ is prime, then $(n+2)\mid n!-1$.
\end{lemma}
\noindent Recall \emph{Wilson's theorem}: a number $p>1$ is prime if and
only if $(p-1)!\equiv-1\pmod p$, i.e.\ the product $1\cdot2\cdots(p-1)$
is one less than a multiple of $p$. We use only the forward direction.
\begin{proof}
For (A) put $p=n+1$; Wilson's theorem gives $n!=(p-1)!\equiv-1\pmod p$.
For (B) put $p=n+2$; then $(p-1)!=(n+1)\,n!$ and $n+1\equiv-1\pmod p$,
so $-n!\equiv-1$, i.e.\ $n!\equiv1\pmod p$.
\end{proof}

\begin{corollary}\label{cor:facprime}
If $n!+1>n+1$ is prime, then $n+1$ is composite. In particular the
factorial primes $n!+1$ with $n\ge3$ (for example $n=3,11,27,37,41,\dots$)
occur only when $n+1$ is composite.
\end{corollary}

\subsection{The worked example \texorpdfstring{$\N=27!+28$}{N=27!+28}}

Take $\N=27!+28$, even with $29$ digits. Since $27!\equiv0\pmod r$ for
every prime $r\le27$,
\begin{equation}\label{eq:shift}
\N-q\;=\;27!+(28-q)\;\equiv\;28-q\pmod r\qquad(r\le27\ \text{prime}),
\end{equation}
so $\N-q$ is killed by a small prime $r$ exactly when $r\mid(28-q)$: a
shifted sieve that eliminates the class $q\equiv28\pmod r$.

\begin{proposition}\label{prop:71}
For $\N=27!+28$ every prime $q$ with $2\le q\le67$ has $\N-q$ composite,
while $\N-71$ is prime; hence $\g(\N)=71$.
\end{proposition}
\begin{proof}[Proof by explicit witnesses]
Table~\ref{tab:witness} gives, for each prime $q\le71$, a factor of
$\N-q$. For $q\le23$ the witness divides $28-q$ via~\eqref{eq:shift}.
Two cases below the winner are notable: $q=29$ is killed by $29$ because
$29=27+2$ is prime, so Lemma~\ref{lem:wilson}(B) gives
$29\mid27!-1=\N-29$; and $q=59$ survives the sieve (as $59-28=31$ is
$27$-rough) yet $41\mid27!-31=\N-59$. Finally $71-28=43$ is $27$-rough
and $27!-43$ is prime.
\end{proof}

\begin{table}[h]
\centering\small
\caption{Why every small prime fails as a summand of $\N=27!+28$.}\label{tab:witness}
\begin{tabular}{rlc@{\qquad}rlc}
\toprule
$q$ & $\N-q$ & factor & $q$ & $\N-q$ & factor\\
\midrule
$2$&$27!+26$&$2$ & $31$&$27!-3$&$3$\\
$3$&$27!+25$&$5$ & $37$&$27!-9$&$3$\\
$5$&$27!+23$&$23$ & $41$&$27!-13$&$13$\\
$7$&$27!+21$&$3$ & $43$&$27!-15$&$3$\\
$11$&$27!+17$&$17$ & $47$&$27!-19$&$19$\\
$13$&$27!+15$&$3$ & $53$&$27!-25$&$5$\\
$17$&$27!+11$&$11$ & $59$&$27!-31$&$\mathbf{41}$ (generic)\\
$19$&$27!+9$&$3$ & $61$&$27!-33$&$3$\\
$23$&$27!+5$&$5$ & $67$&$27!-39$&$3$\\
$29$&$27!-1$&$\mathbf{29}$ (Wilson) & $71$&$27!-43$&\textbf{prime}\\
\bottomrule
\end{tabular}
\end{table}

By~\eqref{eq:shift} a prime $q$ can even be \emph{tested} only if $q-28$
is $27$-rough, so the viable candidates $29,59,71,89,\dots$ are sparse,
and $71$ is merely the first with a prime complement. Sharper still:
$27!+1$ \emph{is} prime (Corollary~\ref{cor:facprime}), the natural
bridge $q=27$, but $27$ is composite, so one falls through to $71$. For
indices that are simultaneously prime and factorial-prime indices
($n=3,11,37,41$) the bridge works and $\g(n!+(n+1))=n$. Heuristically,
the first viable offsets in this factorial family lie on the scale of
$n$, corresponding to roughly $\log\N/\log\log\N$; the rest of the
chapter produces much larger values.

\section{Primorial deserts, made unconditional}\label{sec:primorial}

The factorial wastes each small prime on a residue class. Taking $\N$ to
be a multiple of a primorial wastes nothing.

\begin{lemma}[Primorial core]\label{lem:primcore}
Let $P$ be prime, let $m_P=\prod_{p\le P}p$ be the \emph{primorial} ---
the product of every prime up to $P$ (for instance
$m_{11}=2\cdot3\cdot5\cdot7\cdot11=2310$) --- and let $\N$ be a multiple
of $m_P$ with $\N>2P$. Then $\N-q$ is composite for every prime
$q\le P$; i.e.\ no prime $\le P$ is a Goldbach summand of $\N$.
\end{lemma}
\begin{proof}
For a prime $q\le P$ we have $q\mid m_P\mid\N$ (as $q$ is one of the
factors of $m_P$), hence $q\mid\N-q$, and $\N-q>P\ge q$, so $\N-q$ is a
proper multiple of $q$ and therefore composite.
\end{proof}

Lemma~\ref{lem:primcore} alone gives only claim~(a) of the standard of
evidence. To reach unboundedness with an \emph{actual} Goldbach number
we add Dirichlet's theorem.

\begin{theorem}[Unconditional unboundedness]\label{thm:dirichlet}
For every prime $P$ put $s=\mathrm{nextprime}(P)$. There are infinitely
many primes $t\equiv-s\pmod{m_P}$, and for any such $t$ the even integer
$\N=s+t$ satisfies $m_P\mid\N$, has the explicit Goldbach partition
$\N=s+t$, and has
\[
\g(\N)\;=\;\mathrm{nextprime}(P)\;>\;P.
\]
Hence $\sup_\N\g(\N)=\infty$ unconditionally. Moreover the least such
$t$ satisfies $\log\N\ll P$ by Linnik's theorem~\cite{Linnik}, so
$\g(\N)\gg\log \N$ along this family.
\end{theorem}
\begin{proof}
As $\gcd(s,m_P)=1$ (because $s>P$), Dirichlet~\cite{Dirichlet} gives
infinitely many primes $t\equiv-s\pmod{m_P}$. Then
$\N=s+t\equiv0\pmod{m_P}$ and $\N$ is even. By Lemma~\ref{lem:primcore}
no prime $\le P$ is a summand; the only primes below
$s=\mathrm{nextprime}(P)$ are those $\le P$, so the least summand is at
least $s$. But $\N-s=t$ is prime, so $s$ is a summand; hence
$\g(\N)=s$. Linnik's theorem bounds the least such prime by
$t\ll m_P^{L}$ for an absolute constant $L$, and
$\log m_P=\theta(P)\sim P$, so $\log\N\ll P\sim s=\g(\N)$.
\end{proof}

\begin{remark}[Concrete instances]\label{rem:concrete}
For $P=13$ one has $m_{13}=30030$, $s=17$, and the least prime
$t\equiv-17\pmod{30030}$ is $t=30013$, giving
$\N=30030=2\cdot3\cdot5\cdot7\cdot11\cdot13$ with partition
$30030=17+30013$ and $\g(\N)=17$ --- all primality deterministic. A
larger fully proven example is the primorial
$m_{43}=13082761331670030$ itself: its minimal partition is
$m_{43}=89+13082761331669941$ (both proven prime, being $<2^{64}$), so
$\g(m_{43})=89$ unconditionally.
\end{remark}

\begin{remark}[The roughness bonus]\label{rem:rough}
For the bare choice $\N=m_P$, a prime $q>P$ makes $\N-q$ coprime to all
primes $\le P$ --- the complement is \emph{$P$-rough} --- which improves
its chances of being prime, since the small primes that catch most
composites are already excluded. By \emph{Mertens' theorem} (which says
$\prod_{p\le P}(1-1/p)\sim e^{-\gamma}/\log P$: the product of $1-1/p$
over primes $p\le P$ shrinks like $1/\log P$), this sieving effect
multiplies the heuristic primality probability by
$\prod_{p\le P}(1-1/p)^{-1}\sim e^\gamma\log P$. The expected least
surviving summand then lands near $P(1+e^{-\gamma})\approx1.56\,P$.
Either way the primorial pays a modulus $m_P$ with
$\log m_P=\theta(P)\sim P$ --- exponentially large in the target --- and
the next section removes this waste. The roughness bonus itself,
however, is not waste but the engine of everything that follows: it
reappears in Chapter~\ref{ch:costmodel} as the ``boost'' factor of the
cost model.
\end{remark}

\section{Prime-offset congruence covers}\label{sec:covering}

\subsection{The construction}

A primorial spends one prime per summand; a prime-offset congruence
cover reuses each prime on a whole residue class. The one structural
constraint is parity: a sum of two odd primes is even, so $\N$ is kept
even and the prime $2$ kills only $q=2$; every \emph{odd} prime below
the target must be covered by odd sieving primes.

The strategy itself is classical. Covering systems of congruences were
introduced by Erd\H{o}s precisely to obstruct an additive
representation: his construction of an infinite arithmetic progression
of odd integers not of the form $2^k+p$~\cite{Erdos1950} is the direct
ancestor of this section (and of the Sierpi\'nski--Riesel tradition that
grew from it). What changes here is only the target: the offsets to be
blocked are the primes below $Q$ rather than the powers of two.

\begin{proposition}[Covering certificate]\label{prop:cover}
Let $Q\ge3$ and let $\{(r,a_r):r\in\mathcal R\}$ be a system of odd
primes $\mathcal R$ such that every odd prime $q<Q$ has
$q\equiv a_r\pmod r$ for some $r\in\mathcal R$. If $\N$ is even,
$\N\equiv a_r\pmod r$ for all $r\in\mathcal R$, and
\[
\N\;>\;Q+\max_{r\in\mathcal R}r,
\]
then $\N-q$ is composite for every prime $q<Q$; i.e.\ no prime $<Q$ is a
Goldbach summand of $\N$.
\end{proposition}
\begin{proof}
For $q=2$, $\N-2$ is even and $>2$. For an odd prime $q<Q$, pick
$r\in\mathcal R$ with $q\equiv a_r\equiv\N\pmod r$; then $r\mid\N-q$,
and $\N-q>(Q+\max_r r)-Q\ge r$, so $r\mid\N-q$ properly and $\N-q$ is
composite.
\end{proof}

\begin{remark}
The size hypothesis only has to make every divisor $r\mid\N-q$ proper,
for which $\N>Q+\max_r r$ suffices; the much stronger $\N>Q\prod_r r$ is
unnecessary. A common solution $\N$ exists modulo $M=2\prod_{r}r$ by the
Chinese Remainder Theorem --- the classical fact that congruences to
pairwise coprime moduli can be satisfied simultaneously, and that the
solutions form exactly one residue class modulo the product --- and any
large enough even representative qualifies.
\end{remark}

We build coverings greedily and deterministically.

\medskip
\noindent\fbox{\parbox{0.97\textwidth}{%
\textbf{Algorithm (greedy covering of the odd primes below $Q$).}
\begin{enumerate}\setlength\itemsep{1pt}
\item Let $T=\{\text{odd primes }<Q\}$ and require $\N\equiv0\pmod2$.
\item For $r=3,5,7,11,\dots$ in increasing order, while
$T\ne\varnothing$: let $a_r$ be a residue mod $r$ containing the most
elements of $T$ (smallest such residue on ties), require
$\N\equiv a_r\pmod r$, and delete from $T$ every $q\equiv a_r\pmod r$.
\item Recover $\N$ from the congruences by the Chinese Remainder
Theorem.
\end{enumerate}}}
\medskip

The resulting modulus is $M=2\prod_{r\in\mathcal R}r$ --- twice the
product of the sieving primes actually used. For the computed targets it
is far smaller than the primorial $\prod_{p<Q}p$, because each sieving
prime is reused across an entire residue class rather than assigned to
one offset. (Chapter~\ref{ch:records} replaces this naive greedy pass
with weighted set cover plus annealing; Chapter~\ref{ch:closure} asks,
and largely answers, how far from optimal such covers are.)

\subsection{Partial covers and progression search}\label{sec:partial}

A complete cover is convenient to verify, but it is not the cheapest way
to make $\N$ small. Near the end of the greedy algorithm, a new sieving
prime may remove only one or two remaining offsets while still
multiplying the modulus by $r$. That is an expensive way to buy a single
compositeness proof.

The alternative is to stop early. Let $U$ be the odd prime offsets below
$Q$ not covered by the chosen congruences --- the \emph{residual set}.
CRT gives one residue $\N_0$ modulo
\[
M=2\prod_{r\in\mathcal R}r,
\]
so every candidate has the form
\begin{equation}\label{eq:progression}
\N=\N_0+kM,\qquad k=0,1,2,\dots
\end{equation}
The fixed congruences dispose of most offsets for every $k$. For each
remaining $q\in U$, we test $\N-q$ directly. A failed probable-prime
test is a rigorous compositeness result: pseudoprimes can create false
positives, but they cannot turn a prime into a reported composite. Thus
a fast one-sided test is safe as a search filter, while the final
surviving complement still needs a primality certificate.

\begin{proposition}[Partial-cover certificate]\label{prop:partial}
Let $U$ be the odd primes $q<Q$ not covered by the congruences of
Proposition~\ref{prop:cover}. Let $\N$ be a sufficiently large even
solution of those congruences. If $\N-q$ is composite for every
$q\in U$, then no prime $q<Q$ is a Goldbach summand of $\N$.
\end{proposition}
\begin{proof}
The congruence witnesses prove compositeness for every covered odd
prime; parity handles $q=2$; and the hypothesis handles the residual
set $U$.
\end{proof}

Most values of $k$ fail the hypothesis --- some residual complement
turns out to be prime, which is not a disaster but a small Goldbach
partition, i.e.\ a losing lottery ticket --- and the construction
becomes a \emph{search} along the progression \eqref{eq:progression}.
What that search costs, and how predictably, is the subject of
Chapter~\ref{ch:costmodel}; it is the single question the rest of the
thesis orbits.

\subsection{Realising a covering unconditionally}

Proposition~\ref{prop:cover} still only gives claim~(a). The same
Dirichlet idea as in \S\ref{sec:primorial}, applied twice, upgrades any
covering to an actual Goldbach number.

\begin{theorem}[Unconditional realisation]\label{thm:coverdirichlet}
Let $Q\ge3$ and let $\{(r,a_r):r\in\mathcal R\}$ cover the odd primes
below $Q$, where every $r\in\mathcal R$ is an odd prime. Then there are
infinitely many even integers $\N>Q+\max_{r\in\mathcal R}r$ admitting an
explicit Goldbach partition $\N=s+t$ with $\N\equiv a_r\pmod r$ for all
$r\in\mathcal R$. Every such $\N$ has $\g(\N)\ge Q$.
\end{theorem}
\begin{proof}
For each $r\in\mathcal R$ choose
$b_r\in\{1,\dots,r-1\}\setminus\{a_r\}$, and also require
$s\equiv1\pmod2$. The CRT gives a reduced residue class $b$ modulo
$2\prod_r r$. Dirichlet gives infinitely many primes in that class; call
one of them $s$. If $s<Q$, then $s$ is an odd covered prime, so
$s\equiv a_r\pmod r$ for some $r$, contradicting
$s\equiv b_r\not\equiv a_r\pmod r$. Hence $s\ge Q$.

Next require
\[
t\equiv a_r-s\pmod r\quad(r\in\mathcal R),\qquad t\equiv1\pmod2.
\]
Every residue $a_r-s$ is nonzero modulo $r$, so this is a reduced
residue class modulo $2\prod_r r$. Dirichlet gives infinitely many prime
values of $t$; keep those large enough that $s+t>Q+\max_r r$, and put
$\N=s+t$. Then $\N$ is even, $\N\equiv a_r\pmod r$, and $\N=s+t$ is an
explicit Goldbach partition. Proposition~\ref{prop:cover} excludes every
prime below $Q$, so $\g(\N)\ge Q$.
\end{proof}

\begin{remark}[Quantitative form for initial-segment covers]\label{rem:size}
The covering used in Theorem~\ref{thm:lower} below has one modulus for
every prime $p\le x$. In that setting the two residue classes above are
reduced modulo $M=2\prod_{p\le x}p$. Linnik's theorem~\cite{Linnik}
supplies primes $s,t\ll M^L$ for an absolute constant $L$. Their
residues are nonzero modulo every prime at most $x$, and the covering
property forces the odd prime $s$ to satisfy $s\ge Q$. Consequently the
resulting integer satisfies
\[
\log\N\ll\log M=\theta(x)\ll x.
\]
This quantitative statement is used only for those initial-segment
covers; no such bound is asserted for an arbitrary sparse finite cover.
\end{remark}

\section{Transfer from large prime gaps}\label{sec:gaptransfer}

\begin{proposition}\label{prop:transfer}
If $[\N-Q,\N-2]$ contains no prime, then no prime $q\le Q$ is a Goldbach
summand of $\N$. Consequently, if $\N$ admits a Goldbach representation,
then $\g(\N)>Q$.
\end{proposition}
\begin{proof}
Every prime $q\in[2,Q]$ has $\N-q\in[\N-Q,\N-2]$, hence $\N-q$ is
composite.
\end{proof}

So a prime gap of length $\ge Q$ just below $\N$ forces $\g(\N)>Q$:
\emph{prime gaps transfer to least Goldbach summands}. The transfer is,
however, strictly one-directional. Forcing $\g(\N)>Q$ requires $\N-q$
composite only for \emph{prime} offsets $q<Q$, not for every integer in
an interval; the covering of Proposition~\ref{prop:cover} covers just
the primes below $Q$, a far sparser target. Thus ``large least Goldbach
summand'' is a relaxation of ``large prime gap,'' and one should not
expect the two problems to coincide. Chapter~\ref{ch:gaps} measures the
separation on certified examples; here we exploit the easy direction.

Writing $G(X)$ for the largest gap between consecutive primes below
$X$, the work of Ford--Green--Konyagin--Tao and of
Maynard~\cite{FGKT,Maynard}, sharpened in~\cite{FGKMT}, gives
\[
G(X)\;\gg\;\frac{\log X\,\log_2 X\,\log_4 X}{\log_3 X},
\]
improving Rankin~\cite{Rankin}. In words: there are stretches of
consecutive composite integers below $X$ of length at least a constant
times $\log X\cdot\log_2 X\cdot\log_4 X/\log_3 X$. The dominant factor
is $\log X$ (the average prime gap near $X$); the three iterated-log
factors are a slowly growing extra boost, and since each $\log_k$ grows
glacially this is only a little more than $\log X$. Feeding the covering
that underlies these bounds into Theorem~\ref{thm:coverdirichlet} (which
costs only a constant factor in $\log\N$, by Linnik) makes the transfer
unconditional:

\begin{theorem}\label{thm:lower}
There are infinitely many even integers $\N$, each with an explicit
Goldbach representation, such that
\[
\g(\N)\;\gg\;\frac{\log \N\,\log_2 \N\,\log_4 \N}{\log_3 \N}.
\]
\end{theorem}

\begin{proof}
The Erd\H{o}s--Rankin method, in the sharpened form of
Ford--Green--Konyagin--Maynard--Tao, supplies the covering. Precisely,
\cite{FGKMT} shows (in the standard covering formulation of the method,
the quantity $Y(x)$ of \cite{FGKT,FGKMT}) that there is an absolute
constant $c>0$ such that for all large $x$ one may choose residue
classes $a_p\pmod p$, one for each prime $p\le x$, whose union contains
every integer in $[1,y]$, where
\[
y=\Bigl\lfloor c\,\frac{x\,\log x\,\log_3 x}{\log_2 x}\Bigr\rfloor;
\]
the prime-gap bound is deduced from exactly this covering by the Chinese
Remainder Theorem.

\emph{Normalisation.} If $a_2$ is odd, replace every class $a_p$ by
$a_p-1$; the shifted system covers $[0,y-1]\supseteq[1,y-1]$. So, at the
cost of replacing $y$ by $y-1$, we may assume $2\mid a_2$. Then no odd
integer in $[1,y-1]$ is covered by the class modulo $2$, so every odd
integer $m\le y-1$ --- in particular every odd prime $q<Q:=y-1$ ---
satisfies $q\equiv a_r\pmod r$ for some \emph{odd} prime $r\le x$. The
pairs $(r,a_r)$ over the odd primes $r\le x$ therefore form a covering
of the odd primes below $Q$ in the sense of
Proposition~\ref{prop:cover}: it covers far more (all odd integers below
$Q$), but only the primes are needed, and self-covering ($r=q$,
$a_r=0$) is harmless, exactly as in the computed systems of
Chapter~\ref{ch:records}.

\emph{Realisation.} Theorem~\ref{thm:coverdirichlet} applied to this
covering yields infinitely many even $\N$ with an explicit partition and
$\g(\N)\ge Q$; by Remark~\ref{rem:size} we may take
\[
\log\N\;\ll\;\log2+\theta(x)\;\ll\;x.
\]

\emph{Conclusion.} Let $C$ be the absolute constant with
$\log\N\le Cx$. For large $u$ the function
$u\mapsto u\log u\,\log_3u/\log_2u$ is increasing, and each iterated
logarithm of $\log\N/C$ is $(1+o(1))$ times the corresponding iterated
logarithm of $\log\N$; hence
\begin{align*}
\g(\N)\;\ge\;Q\;&\gg\;\frac{x\,\log x\,\log_3x}{\log_2x}
\;\ge\;\frac{\tfrac{\log\N}{C}\,\log\bigl(\tfrac{\log\N}{C}\bigr)\,\log_3\bigl(\tfrac{\log\N}{C}\bigr)}{\log_2\bigl(\tfrac{\log\N}{C}\bigr)}\\[2pt]
&\gg\;\frac{\log\N\,\log_2\N\,\log_4\N}{\log_3\N}.
\end{align*}
For each sufficiently large $x$, select one realisation $\N_x$
satisfying this Linnik bound. Since $\g(\N_x)\ge Q(x)\to\infty$, the
selected integers are unbounded and therefore contain infinitely many
distinct values. This proves the claim.
\end{proof}

Together with the Granville--van de Lune--te Riele
conjecture~\eqref{eq:GLR}, this brackets the maximal order of $\g$
between a logarithmic-type lower bound and a $(\log\N)^2$-type
conjectural ceiling:
\[
\underbrace{\frac{\log \N\,\log_2 \N\,\log_4 \N}{\log_3 \N}}_{\text{rigorous lower-bound scale}}
\;\ll\;\max_{\N'\le \N}\g(\N')
\;\lesssim\;
\underbrace{(\log \N)^2\log\log \N}_{\text{conjectural ceiling}}.
\]
The explicit constructions of this thesis do not settle that gap; their
value is that they turn the problem into auditable objects --- a
congruence list, a CRT progression, a finite residual set, and a
primality certificate --- and, in Chapters~\ref{ch:limits}
and~\ref{ch:closure}, into measured statements about which parts of the
gap are closable at all. We note one reappearance of this asymptotic
machinery: the FGKMT covering enters again in
Chapter~\ref{ch:closure} not as a lower bound but as the one known
\emph{rounding technology} that might beat greedy covers, which is this
thesis's central open problem.

\section{The first-generation records}\label{sec:firstgen}

The machinery above produced the first constructive records in this
problem family, which the rest of the thesis treats as its baseline.
All three claims (a)/(b)/(c) were discharged to the standard of
Chapter~\ref{ch:intro}, and all were re-verified inside the thesis
pipeline before any successor record was accepted.

\paragraph{The complete-cover baseline.} A complete greedy cover of the
odd primes below $Q=10^5$ uses $356$ congruences and forces a modulus of
$1020$ digits; its least even representative has $\g(\N)=100\,747$.
Every excluded offset has a small factor visible directly from the
congruence system, so verification is little more than modular
arithmetic --- a \emph{proof-oriented build} that overpays roughly
sevenfold in digits for certainty it does not need. (A complete cover at
$Q=10^8$ was also computed as a scaling probe: $327\,455$ digits, with
no partition exhibited.)

\paragraph{The hybrid ablation.} Stopping the cover early and searching
the progression \eqref{eq:progression} shrinks the integer dramatically.
The engineering path, at fixed target $Q=10^5$:

\begin{table}[h]
\centering\small
\caption{First-generation ablation at target $Q=10^5$. ``Congr.''
includes parity; ``residual'' counts prime offsets below $10^5$ left
for direct testing; the last column is the first prime offset whose
complement is (certified) prime.}\label{tab:ablation}
\begin{tabular}{lrrrrrr}
\toprule
Method & Congr. & Digits $M$ & Residual & $k$ & Digits $\N$ & $\g(\N)$\\
\midrule
Complete greedy cover & 356 & 1020 & 0 & 0 & 1020 & $100\,747$\\
Prefix partial cover & 168 & 416 & 385 & 74 & 418 & $101\,599$\\
Coordinate descent & 110 & 248 & 580 & $1\,751$ & 251 & $100\,981$\\
Sparse weighted cover & 103 & 232 & 610 & $299\,948$ & 237 & $109\,621$\\
\bottomrule
\end{tabular}
\end{table}

\paragraph{The $237$-digit height record.} The final row is the
first-generation height incumbent: a $237$-digit even integer
$\N_{237}=\N_0+299\,948\,M$ over a partial cover of $102$ odd
congruences plus parity ($M$ of $232$ digits, largest modulus $1217$),
with
\[
\g(\N_{237})=109\,621.
\]
Every prime offset below $109\,621$ is excluded --- $9\,691$ by explicit
congruence divisors and $731$ by strong base-2 compositeness witnesses
--- and the complement $\N_{237}-109\,621$ is ECPP-certified prime.

\paragraph{The $199$-digit dual incumbent.} Re-optimising the cover for
the strict budget $\N<10^{199}$ gives a $199$-digit integer
$\N_{199}=\N_0+876\,554\,M$ ($89$ odd congruences plus parity, $M$ of
$193$ digits) with
\[
\g(\N_{199})=105\,667,
\]
proved by $9\,335$ parity-or-congruence witnesses, $744$ strong base-2
witnesses, and a $34$-step elliptic-curve certificate for the
complement. It entered the thesis as the incumbent for \emph{both}
bounded games: $T(100\,000)\le\N_{199}$ and $R(10^{199})\ge105\,667$.

\medskip
These records made the covering paradigm concrete, and they framed the
questions the thesis then had to answer. Why $237$ digits and not $150$?
Is the greedy cover any good? What does a progression search actually
cost, and can that cost be predicted rather than experienced? The next
chapter builds the measuring instruments; Chapter~\ref{ch:records}
spends them.

\chapter{The cost model: what a desert costs, and how predictably}
\label{ch:costmodel}

\begin{chapterabstract}
A partial-cover construction is a lottery: each progression element
survives only if every residual complement happens to be composite. This
chapter turns that lottery into a measured science. The success density
is $e^{-\Ee}$ with $\Ee=|U|\cdot\mathrm{boost}/\ln\N$; we derive the
model, refine it with the exact survivor-count distribution, and
validate it to $\pm0.1$ nats across five scales and $1.2\cdot10^9$
scanned candidates --- the dice are exactly as loaded as predicted, every
time. On top of the model we build the economics: the digit-budget
identity that splits $\ln\N$ into modulus nats and search nats, the
digit--compute exchange rate, the two binding regimes of the bounded
games, the boost-versus-kills equilibrium that sets cover size, and the
frontier curves $\Ee(D)$ and $\Ee(Q)$ that price every record before it
is attempted. Two consequences frame the rest of the thesis: a
million-desert below $200$ digits exists in abundance heuristically yet
costs $10^{85}$ primality tests to exhibit by covers --- a
construction-versus-existence gap of fifty-nine orders of magnitude
beyond cosmological --- and the calibrated truth for $T(100\,000)$ sits
near $10^{53}$, a floor our $150$-digit record approaches only in
logarithm.
\end{chapterabstract}

\section{The lottery}\label{sec:lottery}

Fix a partial cover: odd prime moduli $\mathcal R$ with residues
$a_r$, parity, the modulus $M=2\prod_{r\in\mathcal R}r$, and the
residual set $U$ of uncovered odd primes below the target $Q$. Every
candidate is $\N=\N_0+kM$, and $k$ wins exactly when all $|U|$
residual complements $\N-q$ ($q\in U$) are composite
(Proposition~\ref{prop:partial}). Losing tickets are cheap to recognise:
scan the residual complements and stop at the first probable prime.
Winning tickets are certified as in Chapter~\ref{ch:intro}.

Two structural facts make the lottery analysable. First, for residual
$q$ the complement $\N-q$ is automatically coprime to $2$ and to every
cover modulus: $q\not\equiv a_r=\N\bmod r$ means $r\nmid\N-q$. The cover
therefore pre-sieves its own residuals --- the roughness bonus of
Remark~\ref{rem:rough} in its useful form. Second, the residual
complements at a given $k$ are distinct integers of size $\approx\N$
whose only shared structure is exactly that coprimality (a fact made
precise, and load-bearing, in Chapter~\ref{ch:closure}).

\section{The density model}\label{sec:density}

A random integer near $\N$ is prime with probability about $1/\ln\N$.
Conditioning on coprimality to a set $S$ of primes divides the candidate
space by $\prod_{p\in S}(1-1/p)$, so the conditional primality
probability for one residual complement is
\[
p_1\;=\;\frac{\mathrm{boost}}{\ln\N},
\qquad
\mathrm{boost}\;=\;2\prod_{r\in\mathcal R}\frac{r}{r-1},
\]
the factor $2$ from oddness. When the cover contains every prime up to
$r_{\max}$, Mertens' theorem gives
$\mathrm{boost}\approx2e^{\gamma}\ln r_{\max}$; real covers skip some
primes and the product is computed exactly. Treating the $|U|$ residual
primality events as independent (the whole content of that assumption is
examined in Chapter~\ref{ch:closure}; empirically it holds to
$\le0.1$ nats), the probability that a candidate wins is
\[
\prod_{q\in U}(1-p_1)\;\approx\;e^{-|U|\,p_1}\;=\;e^{-\Ee},
\qquad
\boxed{\;\Ee\;=\;\frac{|U|\cdot\mathrm{boost}}{\ln\N}\;}
\]
and the expected number of progression elements to the first hit is
$e^{\Ee}$. The exponent $\Ee$ is the single number that decided every
campaign in this thesis: it is the negative log of the per-candidate
success density, measured in nats.

\paragraph{The exact refinement.} In practice the search engine sieves
each block of $k$ by all primes up to a depth $B_{\mathrm{sieve}}$
(beyond the cover moduli), so at each $k$ only $a_k\le|U|$ complements
remain \emph{alive}, and a surviving complement is prime with the deeper
conditional probability $\hat p$. The exact per-$k$ density is
$(1-\hat p)^{a_k}$, and the honest exponent is
\[
\Ee_{\mathrm{exact}}\;=\;-\ln\;\mathbb E_k\bigl[(1-\hat p)^{a_k}\bigr],
\]
the expectation over the sieve's survivor-count distribution (measured
mean $\approx316$, standard deviation $\approx14$ for a typical
$|U|=750$ cover at $199$ digits, with $\hat p\approx0.059$). Because the
average of an exponential exceeds the exponential of the average, the
fluctuation in $a_k$ is worth a small bonus, $e^{\hat
p^2\sigma_a^2/2}\approx0.2$--$0.3$ nats; the naive formula undercounts
by exactly that much, and the measured densities below confirm the
exact form. The same survivor counts drive the search-order
optimizations of Chapter~\ref{ch:records}.

\section{Validation: five scales, \texorpdfstring{$1.2\cdot10^9$}{1.2e9} candidates}\label{sec:validation}

The model was validated before it was trusted, at every scale it was
later used, by comparing the predicted exponent against
$\hat\Ee=-\ln(\text{measured survival rate})$ on fully logged scans.

\begin{table}[h]
\centering\small
\caption{Density-model validation. ``Predicted'' is
$\Ee$ (exact form where computed); ``measured'' is $\hat\Ee$ from the
logged scan. The two toy deserts are small constructions built to be
exhaustively scannable.}\label{tab:validation}
\begin{tabular}{lrrr}
\toprule
construction & digits of $\N$ & predicted $\Ee$ & measured $\hat\Ee$\\
\midrule
toy desert & 26 & 12.53 & 12.51\\
toy desert & 62 & 13.52 & 13.57\\
dual-game record cover & 197 & 18.20 & 18.2\\
threshold record cover & 195 & 17.51 & 17.4\\
GPU threshold cover & 150 & 24.91 & (one hit at $\lambda=0.63$)\\
mega-desert cover & 2\,692 & 6.60 & 6.56\\
\bottomrule
\end{tabular}
\end{table}

Across three fully instrumented campaign covers the running discrepancy
$|\hat\Ee-\Ee|$ stayed below $0.1$ nats over $1.2\cdot10^9$ scanned
progression elements, with per-slice survivor rates within $\pm2\%$ of
prediction across more than $500$ logged slices. The $150$-digit row
deserves its asterisk explained: that campaign scanned its entire
$k$-round $[0,7.5\cdot10^{10})$ end to end ($\sim2.1\cdot10^{12}$
Fermat tests) and collected exactly one hit against a Poisson mean of
$0.63$ --- an unremarkable draw, and the model stayed unfalsified at the
largest scale ever pushed. This validation is what licenses every
forecast in the thesis: when we say a rung costs $e^{27.3}$ candidates,
that number has survived contact with $10^9$-scale reality repeatedly.

We flag the honest converse: the validation also \emph{bounds} every
dream of beating the lottery. Any exploitable correlation between
residual complements, any hidden compositeness bias in constructed $\N$,
would show up as $\hat\Ee<\Ee$; over $1.2\cdot10^9$ candidates nothing
beyond $0.1$ nats ever did. Chapter~\ref{ch:closure} explains why not.

\section{The digit-budget identity and the exchange rate}\label{sec:budget}

Take logarithms of $\N\approx kM$:
\[
\underbrace{\ln\N}_{\text{record size (nats)}}
\;=\;
\underbrace{\ln M}_{\text{modulus nats}}
\;+\;
\underbrace{\ln k}_{\text{search nats}}.
\]
Modulus nats are spent at design time on congruences (each modulus $r$
costs $\ln r$ and kills a residue class of offsets); search nats are
spent at run time scanning the progression ($\ln k^*\approx\Ee$ at the
expected hit). The two currencies buy the same digits, and the whole
craft of the bounded games is arbitraging them:

\begin{itemize}
\item \textbf{Search nats cost only compute.} A 4-core CPU pipeline
demonstrated $\ln k\approx20$; the GPU fleet of
Chapter~\ref{ch:records} afforded $\ln k\approx25$. Throughput converts
to record digits at the model's slope --- which is why the single
largest threshold-game jump in this thesis ($179\to150$ digits) was
purely an engineering event: the CUDA engine's $430\times$ throughput is
$\ln430\approx6.1$ nats, precisely the exponent gap ($\Ee=24.9$ vs
$\approx19$) between the two rungs.
\item \textbf{The exchange rate is the frontier slope.} Near $200$
digits, $d\Ee/dD\approx-0.10$ to $-0.19$ per digit (so $10\times$
compute buys $12$--$23$ digits there); by $150$ digits the slope is
$\approx-0.32$, and near $100$ digits $\approx-0.5$ ($10\times$ buys
only $\sim5$). The certified ladder measured it directly: each
$\sim7$-digit rung of the $199\to179$ descent cost $\approx2.5\times$
the previous search.
\item \textbf{Cover quality is the third currency.} At fixed $\ln M$,
better residue choices shrink $|U|$; each $1\%$ of $|U|$ is worth
$e^{0.01\cdot\Ee}$ of search. Chapter~\ref{ch:closure} shows the greedy
covers are locally perfect and globally within an integrality gap worth
at most a few nats --- the least solved, most valuable margin in the
game.
\end{itemize}

\section{Two binding regimes, and the equilibrium that sizes covers}\label{sec:regimes}

\paragraph{Compute-bound versus range-bound.} With no size cap, the only
constraint is affordable search: pick the cover minimising
$\ln M+\Ee$ and scan. Under a cap $\N<B$ (the bounded games), the
progression only holds $k_{\max}=B/M$ candidates, and the binding
constraint can flip from compute to \emph{range}. In the range-bound
regime the optimum counterintuitively \emph{shrinks} the modulus below
its unconstrained value: each nat removed from $\ln M$ buys a full nat
of $k$-range while costing only $\approx0.2$ nats of exponent (the
frontier slope in disguise). This was derived, then confirmed twice in
campaign post-mortems: raw GPU throughput pointed at a $200$-digit
budget target without re-optimising $M$ failed to beat a tuned CPU
cover, and the tuned cover's own success confirmed the shrunken-modulus
optimum.

\paragraph{The boost-versus-kills equilibrium.} What stops a cover from
growing forever? Adding a modulus $r$ with its best residue class kills
some number of residuals ($\mathrm{kills}\approx|U|/(r-1)$ on random
structure, less at the picked-over endgame), lowering $\Ee$ by
$\mathrm{kills}\cdot\mathrm{boost}/\ln\N$; but it also multiplies the
boost by $r/(r-1)$, \emph{raising} every surviving residual's primality
odds and hence $\Ee$ by $|U|\cdot\mathrm{boost}/((r-1)\ln\N)$. The two
cancel at
\[
\mathrm{kills}\;\approx\;\frac{|U|}{r-1},
\]
and past that point new moduli are net-zero or worse. Measured at
$B=10^{200}$, $Q=122\,000$: adding $r=461$ to the standing $88$-modulus
cover kills $3.5$ residuals ($-0.084$ nats) and boosts the survivors
($+0.082$ nats) --- net $+0.00$. Two consequences. First, at large
budgets the cover is \emph{not} budget-limited ($\ln M$ used $437$ of
$\sim458$ available nats in that instance); it sits at the equilibrium,
which is why leftover digit budget buys nothing and why the frontier is
flat against re-derivation --- independent greedy runs at the same $Q$
land within $10^{-3}$ nats of each other. Second, the binding constraint
crosses over: at $Q=122\,000$ under $199$ digits the \emph{budget}
binds, while at $Q=200\,000$ the \emph{equilibrium} binds. Both regimes
were mapped before the campaigns that used them.

\section{The frontier curves}\label{sec:frontiers}

Everything above compresses into two curves, computed by building actual
covers (lazy-greedy, digit-capped, annealed) and evaluating $\Ee$
exactly.

\begin{table}[h]
\centering\small
\caption{The digit frontier $\Ee(D)$ at fixed threshold $Q=100\,003$
(greedy digit-capped covers). ``Candidates'' is the expected search
$e^{\Ee}$.}\label{tab:ED}
\begin{tabular}{rrrrr}
\toprule
$D$ (digits) & moduli & $|U|$ & $\Ee$ & candidates\\
\midrule
100 & 52 & 973 & 41.3 & $10^{17.9}$\\
120 & 60 & 899 & 32.8 & $10^{14.2}$\\
130 & 64 & 862 & 29.5 & $10^{12.8}$\\
140 & 68 & 830 & 26.8 & $10^{11.6}$\\
150 & 72 & 802 & 24.4 & $10^{10.6}$\\
160 & 76 & 773 & 22.3 & $10^{9.7}$\\
170 & 79 & 754 & 20.6 & $10^{8.9}$\\
195 & 89 & 692 & 16.8 & $10^{7.3}$\\
\bottomrule
\end{tabular}
\end{table}

\begin{table}[h]
\centering\small
\caption{The threshold frontier $\Ee(Q)$ at fixed budget $199$ digits
(same $88$-modulus digit budget, greedy + coordinate ascent + kicks;
boost $11.01$). The last column prices each rung against its
predecessor.}\label{tab:EQ}
\begin{tabular}{rrrrr}
\toprule
$Q$ & $|U|$ & $\Ee$ & candidates & $\times$ previous\\
\midrule
122\,000 & 867 & 20.83 & $1.1\cdot10^{9}$ & ---\\
135\,000 & 1\,004 & 24.12 & $3.0\cdot10^{10}$ & $\times27$\\
150\,000 & 1\,138 & 27.34 & $7.5\cdot10^{11}$ & $\times25$\\
175\,000 & 1\,335 & 32.07 & $8.5\cdot10^{13}$ & $\times113$\\
200\,000 & 1\,569 & 37.69 & $2.4\cdot10^{16}$ & $\times280$\\
\bottomrule
\end{tabular}
\end{table}

Two planning corollaries. Because rung spacing is geometric, climbing
\emph{every} rung of a ladder costs only $\approx4\%$ more than jumping
straight to the top --- and banks a certified record, plus a fresh
calibration of $\hat\Ee$, at each step. And each curve exposes its
regime: Table~\ref{tab:ED} is the compute-bound threshold game (digits
fall as affordable $\Ee$ rises), Table~\ref{tab:EQ} the range-bound
budget game (the target rises until the equilibrium chokes it).

\section{Construction versus existence: the million-desert study}\label{sec:million}

How far can the paradigm be pushed? As a stress test we priced the
\emph{million desert}: an even $\N<10^{200}$ with $\g(\N)>10^6$, i.e.\
all $78\,498$ primes $q\le10^6$ blocked at once. Note the conjectural
ceiling \eqref{eq:GLR} evaluates to
$(\ln\N)^2\ln\ln\N\approx1.3\cdot10^6$ exactly at $200$ digits: this
target sits \emph{at} the extremal constant of the
Granville--van de Lune--te Riele law, not beyond it. And existence is
not the issue: for unstructured even $\N$ near $10^{200}$ the
all-composite exponent is $\Ee_{\mathrm{random}}\approx445$, so the
expected count of qualifying integers below $10^{200}$ is
$10^{200}e^{-445}\approx10^{6.5}$ --- they exist in abundance,
heuristically.

Exhibiting one is another matter. The frontier $\Ee(D)$ at $Q=10^6$:

\begin{table}[h]
\centering\small
\caption{The million-desert frontier (lazy-greedy digit-capped covers,
$Q=10^6$).}\label{tab:million}
\begin{tabular}{rrrrr}
\toprule
$D$ (digits) & moduli & $|U|$ & $\Ee$ & search cost $e^{\Ee}$\\
\midrule
199 & 91 & 8\,157 & 196.7 & $10^{85.4}$\\
300 & 127 & 7\,311 & 124.3 & $10^{54.0}$\\
400 & 161 & 6\,728 & 89.3 & $10^{38.8}$\\
600 & 226 & 5\,865 & 54.9 & $10^{23.8}$\\
800 & 288 & 5\,241 & 38.2 & $10^{16.6}$\\
1000 & 347 & 4\,741 & 28.4 & $10^{12.3}$\\
1300 & 435 & 4\,100 & 19.6 & $10^{8.5}$\\
1600 & 519 & 3\,588 & 14.2 & $10^{6.2}$\\
2400 & 736 & 2\,492 & 6.9 & $10^{3.0}$\\
\bottomrule
\end{tabular}
\end{table}

The $2400$-digit row reproduces the regime of the certified mega-desert
of Chapter~\ref{ch:gaps} ($\Ee\ 6.9$ vs its measured $6.56$) --- the
ladder is consistent end to end. Reading off feasibility: a
record-scale search ($1.4\cdot10^8$ candidates) needs
$\sim$$1\,470$ digits; a heavy single box ($10^9$) $\sim$$1\,250$; a
distributed fleet ($7\cdot10^{10}$) $\sim$$1\,120$; and a cosmological
budget ($10^{12}$ tests per second for the age of the universe,
$\Ee\le68$) still needs $\sim$$520$ digits. A verified $1\,250$-digit
blueprint ($420$ moduli, $|U|=4\,205$, $\Ee=20.74$, about $235$
core-days) is committed in the artifact repository as the realistic
version of this target.

At $199$ digits the wall is structural. The best $199$-digit system
found --- $91$ moduli, residue-annealed --- leaves $|U|=8\,146$ and
$\Ee=196.4$; annealing recovers only $\sim0.3\%$ of $\Ee$. Structural in
the sieve sense, not the counting sense: the $91$ best classes hold
$148\,290$ slots for $78\,497$ targets, a $1.89\times$ supply, so a
union bound alone forbids nothing --- but prime residue classes overlap
like independent sieves, coverage stalls near $90\%$, and $60\,000$
anneal moves reclaim $11$ residuals out of $8\,157$. Joint optimization
of representative and coverage (Chapter~\ref{ch:limits}) moves
$\Ee$ from $196$ to $\approx194$. So: about $10^{6.5}$ million-deserts
exist below $10^{200}$, and the cheapest known path to \emph{point at
one} costs $10^{85}$ primality tests --- a construction-versus-existence
gap of fifty-nine orders of magnitude beyond a cosmological compute
budget. That gap, in miniature, is also the story of the sub-$100$-digit
threshold game (Chapter~\ref{ch:limits}); the million-desert study is
the thesis's cleanest exhibit that \emph{knowing something exists and
being able to exhibit it are different sciences}.

\section{Poisson discipline: negative results without excuses}\label{sec:poisson}

The lottery is memoryless: conditioned on no hit so far, the expected
remaining wait is unchanged. That single fact, taken seriously, is most
of campaign discipline.

\begin{itemize}
\item \textbf{Spread.} Over the $n=7$ fully-logged searches that ended
in hits, first hits landed at $0.25\times$--$3.0\times$ the expected
$e^{\Ee}$ (median $0.5\times$) on a consistent consumed-candidates
basis. The exponential clock is merciless about variance: plan budgets
at $3\times\Ee$, in writing, before starting.
\item \textbf{Misses are data.} Campaign 1 of the PhD track retired
hitless at $3.18\cdot10^8$ elements --- survival probability $0.30$
under the calibrated model, an unlucky but ordinary draw; per-slice
$\hat p$ and survivor counts matched prediction throughout, ruling out
implementation error. Campaign 2 stopped at $5.6\cdot10^8$ elements
($P\approx0.22$). The model was never adjusted after the fact to
absorb a miss; it never needed to be.
\item \textbf{The tail teaches.} The first-generation-beating
$199$-digit dual record (Chapter~\ref{ch:records}) hit at $99.9\%$ of
its committed range, $3.0\times$ the selection-adjusted expectation ---
the outlier that taught the $3\times$ rule.
\item \textbf{Ratchet, don't grand-slam.} Under fixed compute, a ladder
of certified rungs dominates a single deep campaign: geometric rung
spacing makes the total cost $\approx1.04\times$ the top rung
(\S\ref{sec:frontiers}), each rung banks a record and a calibration,
and optimal stopping becomes arithmetic. The thesis learned this the
expensive way --- a breadth-first 400-variant campaign capped its best
possible outcome and exhausted hitless while a rival single-progression
depth-first design (unbounded runway, bad luck costs digits rather than
termination) banked the record. Both post-mortems are in
Chapter~\ref{ch:records}.
\end{itemize}

\section{Calibrating the truth: how far are the records from $T(Q)$ itself?}\label{sec:truth}

The exhaustive data fix the constant in the conjectural
law~\eqref{eq:GLR}: $\g\le9\,781$ for all even $\N\le4\cdot10^{18}$
gives $C\approx1.4$ in
$\g_{\max}(X)\approx C\,\ln^2\!X\,\ln\ln X$. Inverting at
$Q=10^5$ puts the \emph{true} $T(100\,000)$ near $10^{53}$ ---
plus-or-minus a few in the exponent from the unknown constant and
residue-luck tails --- i.e.\ a $53$-digit integer, unreachable by any
imaginable exhaustive verification (it would take $\sim10^{52}$ scans)
and, Chapter~\ref{ch:limits} will argue, essentially unreachable by
construction too. Our $150$-digit record is thus about $2.9\times$ the
truth \emph{in logarithm}. For the budget game the truth is closer:
$R(10^{200})\approx C\cdot(460.5)^2\ln(460.5)\approx1.8$ million,
against a certified $119\,419$ --- and there the gap is
compute-and-covers, not existence.

For threshold games the model collapses to a useful pair of envelope
identities (cover-quality factor $\varphi\approx0.60$ for current
optimizers; $\ln\mathcal C$ the affordable search exponent):
\begin{align*}
\mathrm{digits}(\N)\;&\gtrsim\;\frac{0.57\,\varphi\,\pi(Q)}{\ln\mathcal C}
&&\text{(compute-bound wall)},\\
\mathrm{digits}(\N)\;&\gtrsim\;\frac{\sqrt{c\,\pi(Q)}}{2.3},\ c\in[0.8,1.3]
&&\text{(absolute floor)},
\end{align*}
the floor being the regime where the ``construction'' has degenerated
into nature's own ensemble scan --- which is why the floor
($38$--$49$ digits at $Q=10^5$) and the truth ($\sim53$ digits) nearly
coincide. Everything between the frontier tables and that floor is the
subject of Chapters~\ref{ch:limits} and~\ref{ch:closure}: how much is
buyable with compute, how much with better covers, and how much is
provably out of reach.

\chapter{Engineering the records: two pipelines, one ladder}
\label{ch:records}

\begin{chapterabstract}
This is the records paper. Two independent pipelines --- a 4-core CPU
stack for cover optimization, search, and certification, and a GPU
fleet on ephemeral cloud notebooks for raw scanning --- drove the
bounded-game records from the first-generation $199$ digits down to
$150$ for the threshold game $T(100\,000)$, and up to
$\g=119\,419$ under $10^{199}$ for the budget game. We describe the
cover optimizer, the search engine and its three measured optimizations
(survivor-ordered testing, $1.8\times$; selective testing, $3\times$
net; sieve-depth tuning), the checkpoint discipline that survived
roughly ten container losses without losing a byte of coverage, and the
five-leg certification standard that survived an adversarial audit ---
including the two occasions a weaker check returned a false PASS and was
caught. Every certified construction is tabulated with its cover
parameters and search cost; the negative campaigns are reported with the
same care, because their statistics are the model validation of
Chapter~\ref{ch:costmodel} in action.
\end{chapterabstract}

\section{Two pipelines}\label{sec:pipelines}

The records were produced by two deliberately independent stacks,
developed in parallel sessions against the same artifact repository.

\begin{itemize}
\item \textbf{The CPU pipeline} runs on 4-core Linux containers:
weighted-set-cover construction with GRASP and simulated annealing,
CRT progression search with a numpy block presieve and gmpy2
strong-probable-prime tests (about $6\cdot10^3$ candidates per second
end to end), and the full certification stack --- PARI/GP ECPP
\cite{PARI,AtkinMorain}, an independent pure-Python certificate
checker, APR-CL~\cite{APRCL} as a second algorithm, and a separate
full-scan crosscheck.
\item \textbf{The GPU pipeline} is a CUDA engine on rotating cloud
notebook sessions: a bit-matrix sieve and Montgomery-CIOS base-2 Fermat
waves, validated bit-for-bit against the Python implementation, at
$2.5$ million tests per second on a T4 and $8.5$ million on an A100 ---
about $430\times$ the 4-core pipeline. A fleet driver shards the
$k$-range across sessions and checkpoints progress in git.
\end{itemize}

The division of labour follows the currencies of
\S\ref{sec:budget}: covers and certificates are latency-tolerant and
optimization-heavy (CPU); the scan is embarrassingly parallel
modular exponentiation (GPU). A shared JSON spec schema (cover pairs,
$\N_0$, $M$, $k$-ranges) is the handoff interface, and every hit found
by either pipeline is certified by both verification stacks before it is
called a record.

\section{Building covers}\label{sec:covers}

The greedy pass of Chapter~\ref{ch:foundations} is the seed, upgraded
three ways.

\paragraph{Weighted set cover.} Choosing modulus $r$ costs $\ln r$ nats
of modulus budget; a candidate pair $(r,a_r)$ is scored by newly covered
offsets per nat. Lazy-greedy selection by this ratio, rather than raw
coverage, is what turned the $1020$-digit complete cover into the
$232$-digit partial cover of the first generation.

\paragraph{Local refinement.} A coordinate-descent pass revisits each
chosen modulus and re-picks its residue against the current uncovered
set, to a fixpoint; GRASP-style randomized ruin-and-recreate with
simulated-annealing acceptance explores further. The refined objective
is $\ln M+\Ee$ --- modulus nats plus search nats, the two currencies on
one scale. Two robust findings: annealing from multiple seeds reliably
reproduces the greedy-plus-descent cover (zero improvement, a plateau
that Chapter~\ref{ch:closure} explains and largely certifies), and the
frontier is flat --- independent re-derivations at the same $Q$ and
budget match $\Ee$ to three decimals.

\paragraph{Variant catalogs.} A single progression under a size cap
holds only $k_{\max}=B/M$ candidates. Residue-swap \emph{variants} ---
same moduli, alternative residues of near-equal coverage --- give
independent progressions at $\approx$zero design cost. The plateau
implies many exactly-equal assignments exist: measured on the
$Q=100\,003$ threshold cover, $26$ zero-penalty single swaps, $325$
pairs, and roughly $2\,600$ triples, a catalog of $\sim3\,000$
progressions at exactly the base $|U|$. Catalogs are what made the
range-bound campaigns possible at all.

\paragraph{An adversarial aside.} An external evolutionary-search
service (AlphaEvolve) was pointed at the cover-step optimizer twice
before this thesis's own closure work adopted it properly
(Chapter~\ref{ch:closure}). The first run's fitness was unconstrained
$-\mathrm{digits}(M)$ from an empty cover, and the evolved program
duly collapsed to infeasible covers ($M$ of $50.7$ digits at
$\Ee=111$) --- a textbook reward hack. With the $\Ee$-cap restored in
the fitness, the evolved move-set improved our best $168$-digit
blueprint by $0.9$ digits at equal wall time; under warm-start
conditions its Metropolis engine lost to our deterministic best-residue
proposals. Verdict: harness patterns adopted, engine not --- and the
incident seeded the anti-reward-hack protocol that
Chapter~\ref{ch:closure} builds on.

\section{The scan engine, and three measured optimizations}\label{sec:engine}

The inner loop is fixed by the mathematics: for each $k$, sieve the
residual complements by small primes, then base-2 Fermat-test the
survivors, abandoning $k$ at the first probable prime. A candidate
costs about $1/\hat p\approx17$ modular exponentiations regardless of
survivor count, because $k$ dies at its first PRP pass; winners are
re-verified exhaustively with BPSW over \emph{all} primes below $Q$
before being reported. Three optimizations, each measured before
deployment:

\begin{enumerate}
\item \textbf{Sieve depth.} Deeper sieving cuts tests per $k$ like
$1/\ln B_{\mathrm{sieve}}$ but costs linearly; the sweep found
$B_{\mathrm{sieve}}=10^6$ optimal on CPU ($901$ vs $816$ candidates/s
at the old $3\cdot10^6$ default), with the measured density invariant
across depths, as it must be.
\item \textbf{Survivor-ordered testing.} Success probability at $k$ is
$(1-\hat p)^{a_k}$: hits concentrate exponentially in the low-$a_k$
tail. Testing each block's $k$ in ascending survivor count, with
stop-at-first-success, gave a measured $1.83\times$ conditional
discovery gain at campaign scale ($12\times$ on a high-density toy;
$53\%$ of toy hits in the bottom survivor-decile). The survivor count
is a sufficient statistic for everything the sieve knows, a fact that
returns with force in Chapter~\ref{ch:closure}.
\item \textbf{Selective testing.} Since candidates are cheap to
generate and expensive to test, skip the fullest ones entirely: with
skip fraction $f=0.92$, raw throughput rose $9.2\times$
($1\,772\to14\,957$ candidates/s) while retaining $29.6\%$ of the hit
mass --- computed \emph{exactly} from the survivor histogram, so the
density model stays unbiased under skipping --- for a net $3.0\times$
discovery rate. The measured retained mass matched the histogram
prediction to $\pm0.4$ percentage points at every $f$ tried.
\end{enumerate}

On GPU the same pipeline becomes a bit-matrix sieve feeding
Montgomery-CIOS Fermat waves; porting optimizations 2--3 to the fleet
was identified as the cheapest remaining nat in the programme.

\section{Fleet discipline: compute that survives its container}\label{sec:fleet}

Both pipelines ran on infrastructure that dies routinely: idle
containers are reclaimed, cloud notebook sessions are culled, processes
vanish mid-slice. The discipline that made multi-day campaigns possible
on such substrate:

\begin{itemize}
\item \textbf{Slices as transactions.} The $k$-space is cut into slices
(e.g.\ $262\,144$ values, $\sim$2 minutes each) interleaved round-robin
across catalog variants; a slice either completes and is checkpointed,
or is re-run. Interleaving makes the global scan breadth-first in $k$,
which also biases the first hit toward small $k$ --- exactly the
threshold game's preference.
\item \textbf{Git as the checkpoint store.} Search logs with resume
marks (\texttt{safe\_k}, per-slice records) are \emph{committed} before
anything else happens. The GPU campaign survived roughly ten session
culls and several container rollbacks with zero coverage loss; the CPU
campaigns absorbed $\sim$15 restarts the same way. Detached processes
and even harness-tracked background tasks measurably did not survive
(duty cycle $12\%$); timeboxed foreground chunks with committed
checkpoints ran at $\sim$85\%.
\item \textbf{Verification at the edge.} Hits are re-verified on the
worker before being reported, so a lost session cannot lose a record ---
harvest first, then die.
\end{itemize}

\section{The certification standard}\label{sec:certification}

Every record passes five legs, uniform across pipelines:
\begin{enumerate}
\item \textbf{Reconstruction}: $\N$ is rebuilt by CRT from the
published congruences and search index and compared byte-for-byte
against the published decimal.
\item \textbf{Exclusion}: every prime $p<\g$ is re-proven a non-summand
from scratch --- proper congruence or trial divisor, else a failed
strong base-2 test --- with the \emph{count} of scanned primes asserted
against an independently computed $\pi(\g)$.
\item \textbf{Partition}: $\g$ is proven prime directly; $\N-\g$'s ECPP
certificate is validated by PARI (\texttt{primecertisvalid}) and
re-verified by an independent implementation
(projective/affine elliptic-curve arithmetic with explicit gcd
surfacing, exact $q>(n^{1/4}+1)^2$ bound checks, chain closed at
$2^{64}$ by deterministic Miller--Rabin), including the binding of the
certificate's top integer to $\N-\g$. CPU-pipeline records add APR-CL
as an algorithmically distinct second proof and a PARI
\texttt{ispseudoprime} full scan of every $p\le\g$ as a third code
path.
\item \textbf{Manifests}: SHA-256 hashes over all artifacts.
\item \textbf{Audit}: an independent session re-ran the verifiers and
mounted deliberate corruptions --- a flipped digit in $\N$, a deleted
witness, a swapped certificate candidate, and a colluding
expected-value edit --- each of which the stack detected loudly.
\end{enumerate}

The standard was not decorative. Twice in the programme's history a
check that skipped a leg returned a false PASS and the remaining legs
caught it: a multi-line GP loop that silently parsed to a vacuous scan
(caught by the prime-count assertion), and a thread-stack overflow that
left a 2-byte junk ``certificate'' (caught by type-checking plus the
independent re-verification). The audit also found three
false-\emph{reject} bugs and zero false-accept paths --- the right
asymmetry, inherited from one-sided witnesses: strong-test failures and
gcd-surfacing arithmetic can wrongly reject, but not wrongly accept. No
record that passed all five legs has ever been overturned. Verification
costs minutes per record; Appendix~\ref{app:artifacts} gives the exact
commands.


\section{The certified ladder}\label{sec:ladder}

\begin{table}[t]
\centering\small
\caption{All certified bounded-game constructions of this thesis, in
order of appearance. ``Odd congr.''\ excludes parity; $|U|$ is the
residual count below the cover's target; the search index is $k$ (or
$t$) in $\N=\N_0+kM$. Game column: T = improved the threshold record
when found, R = improved a budget record, (T)/(R) = certified but
superseded at merge time. Pipeline: C = CPU, G = GPU fleet,
1 = first generation (Chapter~\ref{ch:foundations}).}\label{tab:catalog}
\begin{tabular}{rrlrrrrl}
\toprule
digits & $\g(\N)$ & game & odd congr. & digits $M$ & $|U|$ & index $k$/$t$ & pipe\\
\midrule
1020 & $100\,747$ & baseline & 355 & 1020 & 0 & 0 & 1\\
237 & $109\,621$ & height & 102 & 232 & 610 & $299\,948$ & 1\\
199 & $105\,667$ & T,R & 89 & 193 & --- & $876\,554$ & 1\\
\midrule
199 & $110\,917$ & T,R & 88 & 191 & 729 & $299\,581\,384$ & C\\
197 & $107\,719$ & T,(R) & 88 & 191 & 752 & $951\,928$ & C\\
195 & $100\,297$ & (T) & 88 & 191 & 706 & $70\,186$ & C\\
193 & $102\,337$ & T & 86 & 185 & 686 & $27\,994\,707$ & C\\
186 & $109\,357$ & T & 83 & 177 & 709 & $82\,846\,821$ & C\\
179 & $101\,149$ & T & 80 & 170 & 728 & $758\,850\,385$ & C\\
150 & $104\,527$ & \textbf{T} & 68 & 140 & 827 & $62\,147\,038\,260$ & G\\
\midrule
200 & $112\,249$ & (R) & 89 & 195 & 765 & $56\,770$ & C\\
199 & $119\,419$ & \textbf{R} & 88 & 190 & 806 & $145\,605\,823$ & C\\
\bottomrule
\end{tabular}
\end{table}

Table~\ref{tab:catalog} is the full catalogue; the narrative below walks
it. Parallel sessions occasionally produced overlapping rungs --- the
$195$-digit record landed while a sibling session already held $186$
digits --- and the merged board always kept the best; superseded rungs
are retained as certified constructions because their statistics
validate the model and their covers seed later variants.

\paragraph{Two birds, one integer.} The first thesis records were
\emph{dual-game}: a $199$-digit $\N$ with $\g=110\,917$, built to slip
under the first-generation $\N_{199}$ (it does --- by $0.05\%$, which is
what happens when a hit arrives at $99.9\%$ of the committed range,
$3.0\times$ the selection-adjusted expectation; \S\ref{sec:poisson}'s
cautionary tale), and a $197$-digit $\N$ with $\g=107\,719$, found
after $1.4\cdot10^8$ elements across $165$ progressions at measured
$\hat\Ee=18.2$ against predicted $18.20$.

\paragraph{The threshold descent.} A pure threshold cover needs only
$Q=100\,003$, not the dual game's $105\,668$: $8\%$ fewer targets,
$|U|$ from $752$ to $\sim706$, three times less search
(\S\ref{sec:frontiers}). Stacked with survivor-ordered testing and the
sieve optimum, a two-catalog campaign ($1\,590$ progressions,
$1.75\cdot10^8$ sieved elements, $\hat\Ee=17.4$ throughout) delivered a
$195$-digit record with $\g=100\,297$ on a 4-core box in $\sim$30 hours
of scan --- and, as the tail of the same catalogs, a second independent
hit ($\g=101\,611$ at a larger $\N$), closing the campaign at 2 hits
against a cumulative expectation of $2.1$. Meanwhile the sibling
CPU session ratcheted $193\to186\to179$ digits, each rung costing
$\approx2.5\times$ the previous search, exactly the frontier slope. The
$179$-digit rung was a single progression scanned to
$t=7.6\cdot10^8$ --- it survived $\lambda\approx2.5$ expected hits (an
$8\%$ draw) before landing, the depth-first design whose unbounded
runway converts bad luck into digits rather than termination.

\paragraph{The GPU rung.} At $150$ digits the exponent is
$\Ee=24.91$: about $e^{7.4}\approx1\,600\times$ more search per hit
than at $195$. That rung was unreachable by the CPU pipeline and fell
to the CUDA engine: the full round $[0,7.5\cdot10^{10})$ was sharded
across a rotating two-session fleet, scanned end to end
($\sim2.1\cdot10^{12}$ Fermat tests), and produced exactly one hit at
$k=62\,147\,038\,260$ --- $63.8\%$ into the round, cumulative
expectation $0.40$, a central Poisson draw. The record:
\begin{center}
\begin{minipage}{0.94\textwidth}\small\ttfamily
8266689636142221488914017348682490933831701805981351031435344058125034\\
2747393905081255580217674034223624916359752668869952172293093032932050\\
5713336368
\end{minipage}
\end{center}
$\N_{150}\approx8.27\cdot10^{149}$, $\g(\N_{150})=104\,527$: the
smallest known even integer whose least Goldbach summand exceeds
$100\,000$, $49$ digits below the first-generation record. Evidence:
all $9\,977$ prime offsets below $104\,527$ excluded (parity $1$,
congruence divisor $9\,097$, trial divisor $433$, strong base-2
witness $446$), $104\,527$ prime, complement ECPP-certified.

\paragraph{The budget climb.} Under a hard ceiling the game is
range-bound (\S\ref{sec:regimes}): covers are re-tuned to leave
$k$-room, catalogs supply the progressions, and every hit's
\emph{first surviving offset} usually overshoots the cover target by
thousands --- free upside. The $199$-digit budget record,
\begin{center}
\begin{minipage}{0.94\textwidth}\small\ttfamily
2461381565223403757902559202409729839707029814556797195332164743319183\\
6055009475885630448183435105080229622243237603782026136003397685621411\\
00015486692952424689982503380965035816275690874367307430972
\end{minipage}
\end{center}
has $\g=119\,419$ from a $Q=114\,000$ cover ($88$ odd congruences,
$|U|=806$, $t=145\,605\,823$), bounding $R(10^{199})\ge119\,419$ (and
a fortiori $R(10^{200})$). A parallel campaign at $Q=107\,720$
produced a $200$-digit $\N$ with $\g=112\,249$ --- found in the $16\%$
of sieve blocks that overshoot the $10^{199}$ ceiling, a fair draw at
measured $\hat\Ee=18.4$ against predicted $18.3$, now dominated but
kept: its statistics are part of the validation ledger.

\section{Negative results, reported like records}\label{sec:negatives}

Three campaigns ended without a hit, and one audit prevented a
misallocated one; their numbers appear in the validation of
\S\ref{sec:validation} and their lessons in the design rules above.

\begin{itemize}
\item \textbf{Campaign 1} ($Q=107\,720$, dual-game, $200$ variants):
retired strictly obsolete at $1\,213$ of $2\,400$ slices when the
master merge brought stronger records; $3.18\cdot10^8$ elements, zero
hits, survival probability $0.30$ --- unlucky, ordinary, and fully
on-model per-slice.
\item \textbf{Campaign 2} ($Q=100\,003$, $400$ variants capped at
$175$ digits): stopped at $5.6\cdot10^8$ elements, zero hits,
$P\approx0.22$. Post-mortem against the winning $179$-digit
single-progression design: breadth-first variant pools cap the best
possible outcome and can exhaust; depth-first ratchets dominate under
fixed compute. Adopted.
\item \textbf{Campaign 3} ($Q=122\,000$, budget game): stopped early by
coordination --- a sibling fleet was already attacking the same game at
higher $Q$ --- at $380$ of $17\,200$ slices, hit-mass $\lambda=0.077$,
in-model throughout. Independent workers must scan \emph{disjoint}
variants; duplicated $k$-ranges are pure waste. A committed
variant-ledger became protocol.
\item \textbf{The blueprint audit.} An externally contributed
$168$-digit campaign blueprint was audited before adoption: the cover
was valid, but its claimed $\Ee$ was optimistic by a nat
($19.84$ honest vs $20.85$ claimed), its expected hit was $176.5$
digits (not $168$), and its shipped variants carried a ceiling that
left $\sim$$10^3\times$ too few candidates to matter. Re-parameterized,
it became Campaign 2. Moral: audit contributed specs with the same
machinery as records --- the cost model prices claims, not intentions.
\end{itemize}

\section{Reproduction}\label{sec:reproduction}

Every construction in Table~\ref{tab:catalog} is reconstructable as
$\N=\N_0+kM$ from its committed record file, and every witness is
re-derivable by the standalone verifiers --- nothing depends on the
discovery path. Appendix~\ref{app:artifacts} lists the one-command
verification for each record and the SHA-256 manifests. The next
chapter takes the same machinery to the game with no size limit at all,
where the opponent is no longer the search budget but the primality
certificate.

\chapter{Goldbach deserts are not prime gaps: the height game}
\label{ch:gaps}

\begin{chapterabstract}
A prime gap of length $Q$ just below $\N$ forces $\g(\N)>Q$, so the
height game could in principle free-ride on the prime-gap records. This
chapter shows it should not: blocking only the \emph{prime} offsets is
strictly cheaper than blocking an interval, and the separation is now
certified, not heuristic. Two constructions --- $\g=1\,157\,341$ at
$2\,480$ digits and $\g=1\,134\,871$ at $2\,692$ digits, from
independent pipelines --- exceed the longest prime gap ever certified
with proven endpoints ($1\,113\,106$, between primes of $18\,662$
digits) using integers roughly seven times shorter than that gap's own
endpoints. The second construction additionally carries a fully
certified \emph{local} prime gap of just $8\,970$: its Goldbach desert
is $126.5$ times longer than the prime-free interval it sits in, so the
desert owes essentially nothing to an ordinary gap. At these heights the
search is nearly free ($\Ee\approx6.6$) and the binding constraint
becomes the primality certificate for the partition --- making the
height record an index of primality-proving technology rather than of
number-theoretic difficulty --- and we close with the one construction
idea that would remove that wall too.
\end{chapterabstract}

\section{The transfer, and its one-way sign}\label{sec:transfer5}

Proposition~\ref{prop:transfer} sends prime gaps to Goldbach deserts:
if $[\N-Q,\N-2]$ is prime-free then $\g(\N)>Q$ (given any partition at
all). The converse direction is where the money is, and it fails
constructively: a desert needs $\N-q$ composite only for the
$\pi(Q)$ \emph{primes} $q<Q$, a set of density $1/\ln Q$ inside the
interval the gap must clear entirely. Three certified exhibits measure
the separation.

\paragraph{Exhibit 1: a desert that is a gap in disguise.} The
primorial $m_{43}=13\,082\,761\,331\,670\,030$ of
Remark~\ref{rem:concrete} has $\g(m_{43})=89$, and the nearest prime
below it is $m_{43}-89$ itself: this small desert coincides with an
ordinary prime-free interval. At toy scale the two phenomena are
entangled.

\paragraph{Exhibit 2: a desert 57 times its gap.} The first-generation
$237$-digit record ($\g=109\,621$) already separates them: the
committed verifier proves the interval $[\N-1926,\,\N-2]$ prime-free
while $\N-1927$ is a probable prime ($1927$ is composite, so it is no
summand). Ordinary gap information therefore certifies only
$\g>1\,926$; the actual value is $57\times$ larger, entirely on the
back of composite offsets that the desert may ignore.

\paragraph{Exhibit 3: the certified ratio.} The $2\,692$-digit
construction below carries \emph{both} numbers as certificates: desert
$1\,134\,871$, local prime gap $8\,970$. Ratio: $126.5$.

\section{Benchmarks from the gap world}\label{sec:benchmarks}

The longest prime gap known with \emph{proven} endpoints has length
$1\,113\,106$ and sits between primes of $18\,662$ digits; the longest
known with probable-prime endpoints has length
$16\,045\,848$~\cite{GapList}. Two framing facts. First, certifying a
gap is brutally expensive at scale: the endpoints need primality
proofs, and every interior integer a compositeness witness ---
which is why the proven-endpoint record lags the PRP record by an order
of magnitude. Second, gap \emph{merit} (length divided by $\ln$ of the
endpoint) measures how anomalous a gap is: the largest merit ever
observed is $\approx42$, and Cram\'er--Granville-type
heuristics~\cite{GranvilleCramer} cap plausible merits near a few
hundred. Hiding a million-length desert inside a genuine prime gap at
$200$ digits would need merit $\approx2\,180$ --- fantasy. Deserts at
practical sizes \emph{must} come from covering primes, not intervals;
this chapter's constructions are that statement in certified form.

The challenge the height game set itself: beat $1\,113\,106$ ---
a desert longer than any interval ever proven prime-free --- with
integers small enough to certify on a lab box. (The challenge benchmark
was set at $1\,113\,137$, marginally above the record gap; both
constructions clear both numbers.)

\section{The 2\,480-digit record: $\g=1\,157\,341$}\label{sec:g1}

The CPU pipeline's height construction uses $759$ odd congruence
classes plus parity ($M\approx10^{2474.6}$), leaving $|U|=2\,844$
residual primes below the target $Q=1\,113\,137$, and scans the
progression $\N=\N_0+tM$. At $2\,480$ digits the exponent is tiny ---
$p_1|U|\approx7.7$, expected $\approx$$2.2\cdot10^3$ progression
elements, $\approx1\,800$ \emph{tested} elements under selective
testing --- and the hit arrived at $t=14\,544$, the $443$rd tested
candidate. Certification is where the effort lives:

\begin{itemize}
\item every one of the $89\,840$ primes $p<\g$ proven a non-summand:
$86\,714$ by covering divisor, $3\,126$ by failed strong
probable-prime test, all re-derived by the standalone verifier;
\item $\g=1\,157\,341$ prime by trial division; the complement
$\N-\g$ (a $2\,480$-digit integer) proven prime by a $245$-step ECPP
certificate ($\approx$1.7 hours to generate), re-verified from scratch
by the independent pure-Python checker ($7\,036$ s), \emph{plus}
APR-CL as an algorithmically distinct second proof ($7\,409$ s);
\item an independent PARI full scan of all $89\,841$ primes
$p\le\g$: zero violations ($27.4$ minutes).
\end{itemize}

Result: $\g(\N)=1\,157\,341$, exceeding the certified-gap record by
$4\%$ with an integer $7.5\times$ fewer digits than that gap's
endpoints. The desert record and the gap record are now formally
decoupled: ours needed no interval at all.

\section{The 2\,692-digit mega-desert, and its certified local gap}\label{sec:megagap}

The GPU-session pipeline independently built a second height
construction --- $813$ odd congruences plus parity, $2\,692$ digits,
found at $k=14$ with measured $\hat\Ee=6.56$ against predicted $6.6$
--- with
\[
\g(\N)=1\,134\,871,
\]
proven by $88\,239$ offset witnesses (parity $1$, congruence divisor
$85\,476$, trial divisor $895$, strong base-2 $1\,867$) and an
ECPP-certified complement, the $253$-step certificate re-verified by
the independent checker with its binding to $\N-\g$ asserted.

What makes this construction scientifically distinctive is the
\emph{extra} pair of certificates: the nearest primes around $\N$ are
pinned at $\N-7\,329$ and $\N+1\,641$ --- each carrying its own
validated ECPP certificate, with every interior integer deterministically
composite --- so the ordinary prime gap containing $\N$ has length
exactly $8\,970$, fully certified. The desert-to-gap ratio,
\[
\frac{\g(\N)}{\text{local prime gap}}
=\frac{1\,134\,871}{8\,970}=126.5,
\]
is the thesis's cleanest single number for the claim that Goldbach
deserts are not prime gaps: $99.2\%$ of this desert's length lies
\emph{outside} any prime-free interval, bought by covering prime
offsets alone. In digit terms the separation reads: a $2\,692$-digit
integer carries a $1.13$-million desert, while an interval of that
length free of primes is only known (unproven!) around integers of
$\sim$$18\,000$ digits --- the sparse-cover advantage is about
$7\times$ in digits today and widens with the target.

\section{The economics of height: certification is the wall}\label{sec:certwall}

The cost model explains why the height game is played at thousands of
digits. $\Ee=|U|\cdot\mathrm{boost}/\ln\N$ \emph{falls} as $\N$ grows
--- doubling the digits roughly halves the exponent at fixed cover
strength --- so unlike the digit-frugal bounded games, the height game
optimises at whatever size makes the search trivial:
$\Ee\approx6.6$ at $2\,700$ digits means hits arrive by the dozen. What
does not get cheaper is proving the partition:

\begin{itemize}
\item ECPP generation and validation are minutes at $200$ digits,
$\approx$2 hours (plus 2 hours of independent re-verification, plus
2 hours of APR-CL) at $2\,480$--$2\,692$ digits, and scale roughly
quartically in the digit count;
\item at $\sim$4\,000 digits, days; far beyond that, the certificate
--- not the desert --- is the record's binding constraint with present
tools.
\end{itemize}

The height record is therefore an index of primality-proving
technology. Extrapolating with the cover economics: present ECPP
practice supports $\g\approx2.5$ million (the next committed rung);
distributed fastECPP-class certification ($\sim$$50\,000$-digit
proofs) would support $\g\approx5\cdot10^7$; a certification
breakthrough to $10^6$-digit general-form proofs would carry
$\g$ toward $10^9$. The desert machinery itself is already there: even
the absolute-gap milestone --- beating the $16\,045\,848$ PRP-endpoint
gap, i.e.\ covering $\pi(1.6\cdot10^7)\approx1.03$ million prime
offsets with a $\sim$$20\,000$-digit modulus --- is a throughput
problem ($\sim$$0.2$ s per modular exponentiation at that width; a
distributed campaign $10^2$--$10^3\times$ one lab box), not a
conceptual one. The certificate is the only wall.

\section{Removing the wall: pick the provable prime first}\label{sec:twosided}

Every construction in this thesis searches over $\N$ and pays ECPP for
whichever complement $\N-\g$ falls out --- a random $2\,700$-digit
prime with no exploitable structure. Invert the design: \emph{fix the
complement first}, inside a family whose primality is cheap to
certify, and search the other side.

Concretely, take Proth-form candidates $P=j\cdot2^n+1$ with
$j<2^n$: Proth's theorem~\cite{Proth} certifies such a $P$ prime with a
\emph{single} modular exponentiation once a quadratic non-residue
witness is found --- minutes at $10^5$ digits, where general-form ECPP
is hopeless. For fixed $n$, the cover congruences
$\N\equiv a_r\pmod r$ with $\N=P+s$ ($s$ the intended least summand)
translate to congruences on $j$, since $2^n\bmod r$ is computable; the
search over $j$ is again a CRT progression, now on the \emph{certifiable}
side. A hit hands over the partition $\N=s+P$ with $P$'s certificate
essentially free, deleting the height game's only wall
(Pocklington-style variants~\cite{Pocklington} offer the same trade
with partially factored $P-1$). The open questions are the density of
Proth primes inside the cover-constrained $j$-progression and whether
the $n$-side and $s$-side constraints can be tuned jointly; nobody has
run the construction, and we flag it as the single most valuable open
trick in this problem family --- it would let the $\g$-ladder run to
Proth-search scale ($10^5$--$10^6$ digits, $\g\sim10^8$ and beyond)
with today's tools. Chapter~\ref{ch:conclusions} files it among the
open problems with its neighbours.

\chapter{The limits of compression: why sub-100 digits is out of reach}
\label{ch:limits}

\begin{chapterabstract}
The calibrated truth for $T(100\,000)$ sits near $10^{53}$; the
certified record is $150$ digits. This chapter measures how much of
that logarithmic factor of three is closable, and proves --- in the
measured, adversarially tested sense this thesis uses --- that the
answer is: not much. Four attacks on the size of $\N$ are developed to
their ceilings. Joint optimization of the CRT representative buys
exactly its search entropy ($24.3$ of a provable $24.9$ digits on the
record cover) and cannot compete with progression scanning.
Set-valued conditioning buys $1$--$2.8$ nats, its
extreme-value promise collapsing to $H^{0.35}$, and the classic cover
emerges as the \emph{unique} freely-enumerable design. Oversized covers
would make sub-$100$ trivial ($\Ee=10.2$) but finding their small
representatives is modular subset-sum with per-position choices at
density $\approx1.026$ --- the hardest known regime, with the best
generic attacks at $2^{938}$ here --- and the polynomial-time
CRT-decoding escape is closed twice: by a coupling argument (with a
retraction recorded of our own earlier claim to the contrary) and by
experiment, practical lattice reduction failing to penetrate the
Johnson radius even on solution-abundant instances. The only viable
route to sub-$100$ digits is brute force at
$\Ee=44.1$: about $9\cdot10^5$ GPU-years. The chapter closes with the
walls of all three games and what protects each.
\end{chapterabstract}

\section{The question, and the two axes of attack}\label{sec:question}

The threshold game is a compression problem: pack ``every prime below
$10^5$ fails as a summand'' into as few digits as possible. The
digit-budget identity (\S\ref{sec:budget}) says a construction spends
$\ln\N=\ln M+\ln k$; the frontier (Table~\ref{tab:ED}) prices the
paradigm's trade between them. Beating the frontier means one of two
things: making the \emph{modulus} side cheaper (better covers, or
moduli whose representative is smaller than it deserves to be), or
making the \emph{search} side cheaper (conditioning that raises the
success density). This chapter runs both axes to their measured ends;
Chapter~\ref{ch:closure} then closes the third axis --- replacing
congruences altogether --- with a theorem.

\section{Attack 1: small representatives by construction}\label{sec:joint}

A CRT system's least representative $x$ is what the record pays for; the
modulus $M$ is what covers. Could a system with $M$ far \emph{above}
the budget have $x$ far below it --- oversized coverage smuggled under
the ceiling?

\paragraph{No free small representatives.} Over $2\,000$ random
top-residue reassignments of the $197$-digit record's cover, the
deciles of $x/M$ sit at $(0.10,0.25,0.50,0.75,0.90)$ --- exactly
uniform, with the sample minimum at the $1/\text{samples}$ order a
uniform draw predicts. Residue choice carries no structural bias toward
small representatives; only \emph{search entropy} --- enumerating many
assignments and keeping the smallest --- buys them.

\paragraph{Buying with entropy, at the exact exchange rate.} The tool
is a multiple-choice modular meet-in-the-middle in the style of
Wagner's $k$-tree~\cite{Wagner}: split a block of moduli into $2^d$
lists, enumerate the top-$c$ residue classes of each, and merge with
shrinking centered windows until the representative lands in a target
interval. The entropy accounting predicts $2^d$ lists of size $L$
cancel about $(d{+}1)\log_{10}L$ digits; measured on the record cover:
a 2-list split cancelled $8.3$ digits at a cost of $+8$ residuals
($46$ solutions found where entropy predicted $\approx43$); a 4-level
tree cancelled $21.3$ digits ($52$ vs $\approx50$); the deepest run
cancelled $24.3$ digits against its $24.9$-digit entropy ceiling ---
within half a digit, at $\sim$25 minutes and 8\,GB. Three certified
by-products (committed and re-verified end to end) put oversized moduli
under the $10^{197}$ budget:

\begin{table}[h]
\centering\small
\caption{Joint representative/coverage constructions: modulus far above
the budget, representative under it. The $211$-digit system is a
strictly better cover than the record's own ($|U|$ $738$ vs $752$).}
\label{tab:joint}
\begin{tabular}{lrrrrr}
\toprule
system & odd moduli & digits $M$ & digits $x$ & $|U|$ & $\Ee$ at $10^{197}$\\
\midrule
\texttt{m211\_x196} & 95 & 211 & 196 & 738 & 18.09\\
\texttt{m219\_x197} & 98 & 219 & $<197$ & 784 & 19.33\\
\texttt{m221\_x197} & 98 & 221 & $<197$ & 782 & 19.32\\
\bottomrule
\end{tabular}
\end{table}

\paragraph{Why this does not improve any record.} Extend the objective
to expected hits, $\Phi=\Ee-\ln K(x,M;B)$, where
$K=\max(0,\lfloor(B{-}1{-}x)/M\rfloor{+}1)$ counts progression elements
under the budget. A sieved progression scan generates $\sim$$1\,850$
candidates per second at $\Phi=18.20-\ln(1\,371\,832)=4.07$; $k$-list
construction generates $0.1$--$100$ candidates per second of $K=1$
each --- candidates $10^1$--$10^4\times$ more expensive with $\Ee$ no
lower, $\Phi\ge18$. Joint optimization is the right tool exactly when
the modulus \emph{must} exceed the budget, and there it caps at the
search entropy. The cap is brutal at scale: compressing the
$2\,480$-digit height cover to $197$ digits would need
$\sim$$10^{2495}$ enumerated states --- matching the uniform-draw
estimate $10^{197-2480}$ head-on. Attack 1 is closed by measurement:
the exchange rate is exactly entropy, and entropy is unaffordable.

\section{Attack 2: set-valued conditioning, and the enumerability theorem}\label{sec:setvalued}

The classic cover pins $\N\bmod r$ to \emph{one} residue, paying
$\ln r$ nats; conditioning it into a \emph{set} $A_r$ of $k_r$
residues pays only $\ln(r/k_r)$. A Gaussian extreme-value argument
promises coverage gains of order $\sqrt{2H\sum_rP_r}$ for $H$ nats
spread thinly --- orders of magnitude, if attainable. It is not, and the
reasons are structural.

The exact model (survival as a product over independent
$\N\bmod r$, reproducing the validated
$\Ee=|U|\cdot\mathrm{boost}/\ln\N$ to $0.3\%$ at $k_r\equiv1$) gives,
optimised by Lagrangian sweep over every prime modulus below $15\,000$:

\begin{table}[h]
\centering\small
\caption{Set-valued conditioning at $Q=100\,003$: the entire
achievable gain.}\label{tab:setvalued}
\begin{tabular}{rrrr}
\toprule
$D$ (digits) & classic $\Ee$ & set-valued $\Ee$ & gain (nats)\\
\midrule
100 & 43.84 & 41.83 & 2.0\\
120 & 34.25 & 32.67 & 1.6\\
130 & 31.62 & 29.23 & 2.4\\
140 & 29.36 & 26.53 & 2.8\\
150 & 25.55 & 24.47 & 1.1\\
\bottomrule
\end{tabular}
\end{table}

The promised $\sqrt H$ law measures as $H^{0.35}$, for two reasons:
\emph{overlap} (the nominal kill mass $\sum_rP_r\approx25\,000$ chases
only $\pi(Q)=9\,591$ targets, so most extreme-value gain re-covers
covered primes --- the product structure eats independent
contributions) and \emph{Poisson discreteness} (large-modulus classes
hold $0$ or $1$ targets; there is no excess over the mean left to
collect).

And even those nats are not real, because of \emph{enumerability}. The
set-valued optima condition $\sim$200 moduli of product
$\sim10^{450}$: representatives below $10^{100}$ exist, but finding
which residue combination lands there is precisely Attack 1's problem
again. Free enumeration requires the conditioned modulus $M_1\le B$, in
which case the candidate pool is
$\mathrm{pool}=(B/M_1)\cdot\prod_rk_r$, and optimising under that
constraint the multiplicity term is nearly free but \emph{redundant}:
at $140$ digits it contributes $e^{12.4}$ of pool for $0.15$ nats of
$\Ee$, while plain $k$-room already prices candidates at $0.05$ nats of
$\Ee$ per nat --- net gain $\approx0.2$ nats. The structural conclusion
deserves italics: \emph{the classic cover is optimal not by luck but by
construction --- it is the unique design whose admissible set is an
arithmetic progression, i.e.\ free to enumerate; every relaxation that
buys coverage destroys enumerability, and every relaxation that
preserves enumerability buys almost nothing.}

\section{Attack 3: oversized covers, and the subset-sum barrier}\label{sec:subsetsum}

Suppose we ignore enumerability and aim directly at the dream: covers
whose mass far exceeds the budget. At $B=10^{100}$, $Q=10^5$:
\begin{itemize}
\item \textbf{Regime I} (modulus under the bound): the honest optimum
is $44$ moduli, $M=10^{79}$, $|U|=1\,060$, $\Ee=44.1$ ---
$1.4\cdot10^{19}$ candidates.
\item \textbf{Regime II} (oversized): $420$ moduli with $16$ allowed
classes per position give $|U|=165$ and $\Ee=10.2$ --- about
$27\,000$ candidates. Sub-$100$ would be a lunch-break computation.
\end{itemize}
Regime II's design space carries $\approx2\,698$ nats of entropy
against $\ln(M/B)\approx2\,630$: about $e^{57}\approx2^{82}$ valid
(design, $\N$) pairs \emph{exist} below $10^{100}$. But the density of
designs whose representative lands under the bound is $e^{-2630}$, and
finding one is exactly:

\begin{quote}
\emph{minimise $\bigl|\sum_r c_r \bmod M\bigr|$ over per-position
choices $c_r\in C_r$, $|C_r|=L$, target window $[0,B)$}
\end{quote}

--- inhomogeneous modular subset-sum with choices, at density
(design entropy)/$\ln(M/B)\approx\mathbf{1.026}$. This is the worst
possible regime: low-density instances ($d<0.94$) fall to lattice
reduction \cite{LO,CJLOSS}, very high densities to birthday/$k$-tree
merging \cite{Wagner}, and density $\approx1$ is where subset-sum's
presumed hardness concentrates. Concretely here: meet-in-the-middle
$2^{1897}$; Wagner's $k$-tree with the available list sizes $2^{938}$;
both astronomically short of $2^{82}$ planted solutions in a
$2^{3893}$ space.

\paragraph{The decoding escape, closed --- with a retraction.} One
hope remains: treat the residues as a Chinese Remainder code and use
Guruswami--Sahai--Sudan list decoding~\cite{GSS} (in Howgrave-Graham's
lattice formulation~\cite{HG}) to \emph{decode} a small representative
directly. The verified radius: agreement mass $A>\sqrt{L\ln B\cdot P}$
suffices, $P$ the log-mass of the whole modulus pool (numerically
confirmed against the exact finite-$\ell$ optimum to $0.4\%$). Version
1 of our analysis claimed an open feasibility window for this attack.
\textbf{That claim was wrong and is retracted}: it compared the
decoding radius at one agreement size against an existence bound
optimised over a \emph{different} size. Coupling the parameters closes
the window at every pool size --- the barrier in one breath:
\emph{decodable $\Rightarrow$ pool small $\Rightarrow$ subset entropy
too small for a representative to exist; representative exists
$\Rightarrow$ pool large $\Rightarrow$ radius far beyond the cover
mass.} Re-derived independently (Legendre-transform design entropy,
same-mass coupling):

\begin{table}[h]
\centering\small
\caption{The decoding barrier: maximum cover mass for which
representatives exist vs minimum mass the decoder needs, same pool and
$L$. The window never opens; the best case is $28\%$ short.}
\label{tab:barrier}
\begin{tabular}{rrrrrr}
\toprule
pool $R$ & $L$ & pool mass $P$ & $A_{\max}$ (exists) & $A_J$ (decodes) & ratio\\
\midrule
$1\,000$ & 1 & 956 & 339 & 469 & 0.723\\
$60\,000$ & 1 & $59\,816$ & 570 & $3\,711$ & 0.154\\
$60\,000$ & 16 & $59\,816$ & $1\,305$ & $14\,845$ & 0.088\\
$10^6$ & 16 & $998\,483$ & $1\,697$ & $60\,651$ & 0.028\\
\bottomrule
\end{tabular}
\end{table}

The gap widens with pool size and with $L$ ($A_J\sim\sqrt{LP}$ against
a logarithmic existence bound). The retraction's lesson is recorded for
the methodology chapter: \emph{feasibility windows built from two
separately-optimised bounds are worthless; couple the parameters
first.}

\section{Attack 3, empirically: LLL versus the Johnson radius}\label{sec:toycrt}

Proofs bound algorithms that exist; experiments bound hope. Could
practical lattice reduction beat the proven radius on our
\emph{structured, solution-abundant} instances? A controlled study at
toy scale: $61$ positions (odd primes $<300$, $\ln M=276.3$), messages
below $B=10^{12}$, Johnson radius $A_J=\sqrt{\ln B\ln M}=87.38$;
planted instances revealing mass $\rho A_J$ of true residues; LLL
(FLINT, $\delta=0.99$) on the Howgrave-Graham lattice at degrees
$\ell\in\{12,20,30,45\}$; $390$ decodes total, predeclared protocol.

Results. Recovery is $0\%$ at $\rho\le0.90$, $60\%$ at $\rho=1.00$,
$100\%$ at $\rho\ge1.10$; measured against \emph{realized} agreement
mass the transition is razor sharp: at $\ell=30$ every decode below
$1.0271\,A_J$ failed and every one above $1.0278\,A_J$ succeeded ---
threshold $\approx1.028\,A_J$, tightening to $\approx1.02\,A_J$ at
$\ell=45$, tracking the finite-dimension determinant bound
\[
A_{\min}(\ell)=\min_z\frac{z(z{+}1)\ln M+\ell(\ell{+}1)\ln B}{2(\ell{+}1)z}
\]
to about $1\%$ below it (the usual Coppersmith-family grace) and
converging to $A_J$ \emph{from above} like $1+c/\ell$. Below threshold
the failure is total: the reduced bases yield integer-rootless
polynomials --- not the planted message, and not one of the
$10^2$--$10^7$ high-agreement alternatives that the first-moment count
says exist at $\rho=0.7$--$1.0$. Across all $390$ decodes, every
integer root ever extracted was the planted message itself.
\emph{Solution abundance buys nothing.} The Johnson radius behaves as a
hard practical barrier for this lattice family, closing the last
``maybe practice beats the proof'' hope from the barrier analysis.

\section{The sub-100 verdict, priced}\label{sec:sub100}

With Attacks 1--3 closed, sub-$100$-digit $T(100\,000)$ has exactly one
route: Regime-I brute force at $\Ee=44.1$. Two independent costings
(this thesis's, and the parallel realizable-frontier study, which
lands at $\Ee=43.66$ --- agreement to $0.4$ nats) put it at
$1.4\cdot10^{19}$ candidates $\approx9\cdot10^5$ GPU-years at the
record engine's measured rate --- roughly seven orders of magnitude
beyond plausible compute, and the number to quote against anyone's
``surely a bigger machine\dots''. The full menu, cap-aware (these
$\Ee$ include the $k$-room overhead, hence sit about a nat above the
raw frontier of Table~\ref{tab:ED}; the model reproduces the actual
$150$-digit campaign --- two GPUs, six days --- as its calibration row):

\begin{table}[h]
\centering\small
\caption{The descent below 150 digits, priced. Wall-clock rows from the
measured fleet band ($7.9\cdot10^5$ candidates/s end to end); multiply
by $\sim$1.6 for $80\%$ confidence.}\label{tab:menu}
\begin{tabular}{rrrrl}
\toprule
$D$ & $\Ee$ & candidates & A100-years & fleet wall-clock\\
\midrule
150 & 25.5 & $10^{11.1}$ & 0.01 & $\sim$2 days (done: the record)\\
140 & 28.1 & $1.6\cdot10^{12}$ & 0.10 & $\sim$13 days\\
135 & 29.7 & $7.7\cdot10^{12}$ & 0.49 & $\sim$47 days\\
130 & 31.2 & $3.6\cdot10^{13}$ & 2.3 & $\sim$5 months\\
125 & 32.9 & $1.9\cdot10^{14}$ & 12 & ---\\
120 & 34.7 & $1.2\cdot10^{15}$ & 73 & ---\\
100 & 44.1 & $1.4\cdot10^{19}$ & $8.9\cdot10^5$ & ---\\
\bottomrule
\end{tabular}
\end{table}

So the fundable frontier is $130$--$140$ digits; $120$--$125$ is the
institutional edge ($\sim$$10^2$ GPU-years); below that the record
stalls on physics unless mathematics moves. Not impossibility ---
about $10^{53}$ qualifying integers exist below $10^{100}$
(\S\ref{sec:million}'s accounting at this scale: the unstructured
exponent is $\approx110$ against $230$ nats of room) --- but
constructive inaccessibility, protected on two independent axes:
search entropy in Regime I, and density-$1$ subset-sum in Regime II.
For completeness, the quantum row: Grover~\cite{Grover} halves the
exponent, but with resource estimates in the
Gidney--Eker\aa\ mold~\cite{GE} a fault-tolerant oracle call
(reversible $660$-bit modular exponentiations, no early exit) costs
$\sim$$10^5$--$10^6$ seconds against the classical
$1.25\cdot10^{-6}$: the crossover sits at $\Ee^*\approx55$ for one
machine, $\approx41$ for a thousand --- at or above every target we
care about, with both sides infeasible at the crossover. The walls of
this chapter are classical \emph{and} quantum walls at any constants
that have been projected.

Two well-posed problems would collapse the verdict, and are of
independent interest: \emph{CRT list-recovery beyond the Johnson
radius for structured, solution-abundant instances}, and
\emph{modular subset-sum with per-position choices at density
$\approx1$ when $2^{82}$ solutions exist}. They are filed in
Chapter~\ref{ch:conclusions}.

\section{The walls, per game}\label{sec:walls}

Assembling this chapter with the calibrated truths of
\S\ref{sec:truth} gives each game its forecast. The table uses
$\ln\mathcal C\approx25$ for today's fleet, $\approx35$ for a
nation-scale scan, $\approx45$ as a generous classical ceiling.

\begin{table}[h]
\centering\small
\caption{THE WALL: predicted limits per game. ``Truth'' is the
calibrated Granville--van de Lune--te Riele estimate; the middle
columns are what the validated machinery reaches at stated compute;
``protected by'' names the obstruction that survives everything this
thesis threw at it.}\label{tab:walls}
\begin{tabular}{p{0.16\textwidth}p{0.13\textwidth}p{0.17\textwidth}p{0.16\textwidth}p{0.25\textwidth}}
\toprule
game & truth & today / heroic & likely future & protected by\\
\midrule
$T(10^5)$, smallest $\N$ & $\sim10^{53}$ (53 digits) & $\sim$140d (days), $\sim$130d (months) / 120--125d & stall at 120--135d; floor 38--49d & search entropy; density-1 subset-sum; the closure triple (Ch.~\ref{ch:closure})\\
$R(10^{200})$, largest $\g$ & $\sim$1.8M & $\sim$170k / $\sim$250k & $\sim$450k & range entropy under the cap\\
Height, largest $\g$ & unbounded & $\sim$2.5M / $\sim5\cdot10^7$ (fastECPP) & $10^9+$; unbounded if two-sided construction lands & certification cost only\\
\bottomrule
\end{tabular}
\end{table}

The budget game is where algorithms matter most per unit compute (its
truth is only $\sim$$15\times$ away); the threshold game is where they
matter least (its truth is entropy-armoured from two sides); and the
height game has no mathematical wall at all, only an engineering one.
One axis remains unexamined by everything so far: perhaps congruence
divisors are simply the wrong \emph{mechanism}, and some algebraic
structure certifies compositeness wholesale without paying $\ln r$ per
class. Closing that axis --- or being handed a revolution by failing to
--- is the next chapter.

\chapter{The Divisor-Paradigm Closure Theorem}
\label{ch:closure}

\begin{chapterabstract}
Everything before this chapter certifies compositeness one way: a
planted congruence divisor. Could some algebraic mechanism --- norm
forms, cyclotomic identities, exponent families, polynomial
factorizations --- certify many complements composite wholesale,
escaping the $\ln r$-per-class accounting that walls in
Chapter~\ref{ch:limits}? This chapter answers with a theorem and its
falsification engine. We formalise construction schemes as
candidate families with certificate maps, prove that certificates
either \emph{plant} their divisor at design time or must
\emph{extract} it per instance (which is factoring), and show --- across
a concrete grammar $\mathcal L$ of planted mechanisms --- that
congruence covering is the unique economical one: value-set ceilings
kill univariate identities at any price, a first-moment cluster bound
caps every shared-variable identity family at $O(1)$ certified offsets,
isolation and orbit accounting close the window-dense frontier, and a
supply-side purchase argument closes the cheapest surviving family
$2.5\times$ over budget. The engine swept $13\,661$ enumerated and $41$
adversarially evolved instances: zero passes --- but the adversary's
first find was a $50\times$ accounting hole in our own referee, and
fixing it sharpened the theorem's economics. With mechanisms closed,
what governs the records is measured too: the entire game exponent is
an integrality gap (the fractional cover is perfect), and the greedy
point survives exact optimization over windows up to half the cover.
Three doors remain --- rounding mathematics, throughput, and two named
conditionals --- and we state each at its maximal form.
\end{chapterabstract}

\noindent\emph{Rigor marks, used throughout this chapter:} [P] proved at
this thesis's level of rigor; [C] conjectural; [V] empirically
validated input with the dataset stated. The theorem's honest status is
assembled from these marks in \S\ref{sec:verdict}.

\section{Five probes before the theorem}\label{sec:probes}

Five bounded probes, each with a predeclared termination gate, tested
the folk candidates for a better mechanism. The benchmark: covering
certifies $\sim$$4\,400$ offsets per nat of design entropy at $r=3$,
falling to $\sim$$1$ per nat at the endgame ($r\approx457$; measured
marginal $1.01$ nats per offset averaged over the last six moduli).

\begin{itemize}
\item[\textbf{P1}] \emph{Prime-power and composite moduli.} Terminated,
strictly dominated: the class $q\equiv a\ (\mathrm{mod}\ p^2)$ is a
subset of the class mod $p$ already killed, at double the cost
[V: $p=7$ gives $227\subset1\,613$ kills]. Composite moduli reduce to
their prime factors with only losses.
\item[\textbf{P2}] \emph{Polynomial value-set identities.} $\N=m^k-c$
makes $\N-q=m^k-d^k$ algebraically composite for every offset
$q=d^k-c$. Terminated, dominated --- but the closest call, and the
seed of Lemma~1: the best quadratic family kills $113$ offsets below
$10^5$ [V: $c=398$], yet confining $\N$ to a $k$-th-power family
surrenders $(1-1/k)\ln B$ nats, and \emph{the value set of a
degree-$k$ form below $Q$ has $O(Q^{1/k})$ members} --- at most a few
hundred killable offsets at any price.
\item[\textbf{P3}] \emph{Sierpi\'nski--Riesel exponent families.}
$\N_n=j\cdot2^n+c$: for any fixed candidate the induced constraint is
one congruence class --- ordinary covering. The exponent structure
correlates residues \emph{across} the family, a legitimate
search-throughput trick, not an entropy mechanism. Terminated.
\item[\textbf{P4}] \emph{Norm forms and splitting conditions.} ``$\N-q$
is a norm in $K$'' does not certify compositeness (primes with the
right splitting are norms too), and any usable splitting condition on
$q$ is a congruence condition mod discriminant-related moduli ---
covering in Galois clothing. Terminated.
\item[\textbf{P5}] \emph{The hidden-bias loophole.} Could constructed
$\N$ make residual complements composite more often than the boost
predicts? The banked campaigns answer: across three covers and
$1.2\cdot10^9$ scanned candidates, $|\hat\Ee-\Ee|\le0.1$ nats [V].
Any unexploited structure is worth $\le0.1$ nats. Terminated.
\end{itemize}

The probes sharpen the requirement. A paradigm-breaking mechanism must
\emph{simultaneously} (1) certify $\gg\sqrt Q$ distinct prime offsets,
(2) cost $o(\ln r)$ design entropy per offset at the margin, and (3) be
invisible in banked density data --- activating only on constructed
$\N$. No known structure does all three. The rest of the chapter turns
that requirement into a theorem about a language.

\section{Schemes, certificates, and the theorem}\label{sec:formalism}

Fix game parameters: bound $B$ (here $10^{200}$), threshold $Q$ (here
up to $2\cdot10^5$), $\ln\N\approx\ln B$.

\begin{definition}[Construction scheme]
A \emph{construction scheme} is a pair $(F,C)$:
\begin{itemize}
\item a \emph{candidate family} $F$: a parameter set
$\Theta\subseteq\mathbb Z^m$ with polynomial-size description and an
evaluation map $\N:\Theta\to2\mathbb Z\cap[1,B)$ computable in
$\mathrm{poly}(\log B)$; its \emph{design cost} is
$\mathrm{cost}(F)=\ln B-\ln|\N(\Theta)|$, the nats of search freedom
surrendered by joining the family;
\item a \emph{certificate map} $C$: on input $(\theta,q)$, $q$ prime
$<Q$, outputs $\bot$ or a divisor witness $D(\theta,q)$ with
$1<D<\N(\theta)-q$ and $D\mid\N(\theta)-q$, computable and verifiable
in $\mathrm{poly}(\log B)$.
\end{itemize}
$C$ is \emph{planted} when the divisibility holds as an algebraic
identity over the family (checkable symbolically); \emph{extracted}
when $C$ discovers the divisor per instance. The \emph{killed set} is
$K(\theta)=\{q:C(\theta,q)\ne\bot\}$, the \emph{residual set}
$U(\theta)$ its complement in the primes below $Q$, and the
\emph{effective exponent} $E_{\mathrm{eff}}(F,C)$ is the infimum, over
polynomial-time searches, of the log of expected work per hit divided
by work per candidate. For the classic cover,
$E_{\mathrm{eff}}=\Ee=|U|\cdot\mathrm{boost}/\ln\N$
[V: \S\ref{sec:validation}].
\end{definition}

The \emph{covering frontier} $\Ee^*(Q,B)$ is the minimum $\Ee$ over
congruence covers at the boost-versus-kills equilibrium of
\S\ref{sec:regimes} (e.g.\ $\Ee^*(122\,000,10^{199})=20.83$;
$\Ee^*(200\,000,\sim10^{200})\approx37.7$).

\begin{theorem}[Divisor-Paradigm Closure --- target statement {[C]},
proved on the grammar {[P]}]\label{thm:closure}
Let $(F,C)$ be a construction scheme whose certificate map is planted
and expressible in the schema language $\mathcal L$ of
\S\ref{sec:grammar}. Assume:
\begin{itemize}
\item[\emph{(H1)}] Hardy--Littlewood-type independence for shifted
primes in families: conditional on all divisibility data by primes
$<Q$ (and the sieve depth), the primality indicators of distinct
residual complements are independent with the Mertens-adjusted rates;
\item[\emph{(H2)}] no compositeness test with constant soundness on
random rough inputs runs in $o$(one modular exponentiation).
\end{itemize}
Then $E_{\mathrm{eff}}(F,C)\ge\Ee^*(Q,B)-o(1)$. Moreover, extracted
certificates on any window-dense family of targets are
factoring-equivalent, and certificate-free detection is already
employed at its known-optimal cost (one Fermat exponentiation) by the
standard engine.
\end{theorem}

Informally: \emph{within $\mathcal L$, nothing beats covering; outside
planted-$\mathcal L$, a scheme must either factor or test, and both sit
at known cost floors.} Before the lemmas, a sanity audit of everything
outside the formalism, from the million-desert study
(\S\ref{sec:million}): perfect powers cannot help (consecutive squares
near $10^{200}$ are $\sim10^{100}$ apart, so the window contains none);
hiding the desert in a genuine prime gap needs merit $\approx2\,180$
against $\approx42$ ever observed (\S\ref{sec:benchmarks}); composite
moduli are P1. The formalism is aimed at the only survivors: algebraic
identities and per-instance discovery.

\section{The schema language $\mathcal L$}\label{sec:grammar}

The falsifiable core of the theorem is a concrete grammar over
certificate mechanisms; each compiles to an offset enumerator, a
symbolic divisor witness, and a design-cost term:
\begin{align*}
\mathcal L\ ::=\ &\textsf{Congruence}(d,a) &&\N\equiv a\ (\mathrm{mod}\ d)\\
\mid\ &\textsf{PolyImage}(f\in\mathbb Z[m],\deg\le4) &&\N-q=A(m)\cdot B(m)\ \text{identically}\\
\mid\ &\textsf{ExpFamily}(\N=j\cdot a^n+c) &&\text{per-offset multiplicative-order conditions}\\
\mid\ &\textsf{MultiPoly}(A,B\in\mathbb Z[u_1..u_v],\deg\le2,v\le3) &&\N-q=A(u)\cdot B(u)\\
\mid\ &\textsf{NormForm}(K,\ \mathrm{disc}\le10^4) &&\N-q=\mathrm{Norm}_K(\alpha),\ \text{exhibited factorization}\\
\mid\ &\textsf{Cyclotomic}(\Phi_n,\ n\le200) &&\text{cyclotomic/Aurifeuillian splittings}\\
\mid\ &\textsf{Compose}(\mathcal L,\mathcal L) &&\text{depth}\le2.
\end{align*}
Excluded by construction: ``certificates'' that need a primality
decision to define membership (circular), and per-instance divisor
discovery --- the extraction horn, handled by the theorem's second
clause rather than by the grammar. Grammar extensions are a human act
with a hardening rule (\S\ref{sec:engine7}): any extension to $v\ge4$
shared variables must ship with the dense-frontier analysis of
Lemma~1(c) re-proved first.

\section{Lemma 1: entropy accounting for planted certificates}\label{sec:lemma1}

\emph{Statement.} For planted $(F,C)$, the amortised design cost per
certified offset is at least $\ln(1/\delta_{\mathrm{eff}})-O(\ln\ln B)$,
where $\delta_{\mathrm{eff}}$ is the density of integers in a generic
window of width $Q$ at height $B$ for which the scheme can exhibit its
witness with polynomial-range parameters; and for every family in
$\mathcal L$ this is no cheaper than covering. The proof is a case
sweep, and the cases are the chapter's real content.

\paragraph{(a) Congruence classes [P].} $\delta_{\mathrm{eff}}=1/d$,
cost $\ln d$: the covering accounting itself, with composite $d$
dominated by its prime factors (P1).

\paragraph{(b) Univariate identities [P].} A degree-$k$ identity's
certified offsets lie in the value set of a degree-$k$ form below $Q$:
at most $O(Q^{1/k})$ of them, \emph{at any price}. At $Q=122\,000$:
$349$ ($k{=}2$), $50$ ($k{=}3$), $19$ ($k{=}4$), $7$ ($k{=}6$); and
$\Phi_n$ has degree $\varphi(n)\ge2$ for $n\ge3$, so
\textsf{Cyclotomic} sits inside the same ceiling. Both families are
closed \emph{price-free}: the engine's pass bar (below) demands
$>3\sqrt Q\approx1\,048$ distinct offsets, and $349<1\,048$.

\paragraph{(c) Multivariate identities --- the hard case.} Here the
representable set can be window-dense (products of two binary quadratic
forms reach $\sim$$9\,300$ of the $17\,983$ prime offsets at
$Q=200\,000$ [V]), so existence arguments must be replaced by
\emph{access} arguments. Write the factors' degrees $d_1,d_2\ge2$
(degree-1 factors are sieve-equivalent), $d=d_1+d_2$, $v$ shared
variables, and $\alpha=\min(1,v/d)$ the value-set density exponent
[V: measured $\hat\alpha=0.761$ vs $0.750$ predicted for $(2,2,3)$].

\emph{Cluster ceiling, $v<d$} [P modulo standard equidistribution of
form values]. Distinct values of $G=A\cdot B$ below $X$ number
$\sim X^\alpha$; a first-moment count of $k$-clusters inside width-$Q$
windows gives $X^\alpha(X^{\alpha-1}Q)^{k-1}$, so clusters exist only
up to
\[
k_{\max}(\alpha)=1+\Bigl\lfloor\frac{\alpha\ln X}{(1-\alpha)\ln X-\ln Q}\Bigr\rfloor,
\]
finite iff $\alpha<1-\ln Q/\ln X$ (threshold $\alpha^*=0.9735$ at
$X=10^{200}$, $Q=2\cdot10^5$). Every $v<d$ combination has
$\alpha\le(d-1)/d\le0.9<\alpha^*$:

\begin{table}[h]
\centering\small
\caption{The cluster ceiling: maximum certifiable offsets per candidate
for shared-variable identity families $(d_1,d_2,v)$ at $X=10^{200}$,
$Q=2\cdot10^5$. A scheme needs $\sim$$18\,000$.}\label{tab:kmax}
\begin{tabular}{lrr@{\qquad}lrr}
\toprule
combo & $\alpha$ & $k_{\max}$ & combo & $\alpha$ & $k_{\max}$\\
\midrule
$(2,2,2)$ & 0.500 & 2 & $(2,2,3)$ & 0.750 & 4\\
$(2,3,3)$ & 0.600 & 2 & $(2,3,4)$ & 0.800 & 5\\
$(2,4,4),(3,3,4)$ & 0.667 & 3 & $(2,4,5),(3,3,5)$ & 0.833 & 6\\
$(3,4,5)$ & 0.714 & 3 & any $v<d$ & $\le0.9$ & $\le13$\\
\bottomrule
\end{tabular}
\end{table}

For the flagship case $(2,2,2)$ this is existence, not hardness: at
$X=10^{200}$ the expected number of 3-clusters is
$Q^2X^{-1/2}\approx10^{-89}$, so a same-variable quartic-product scheme
certifies \emph{at most two} offsets per candidate with high
probability over any construction ensemble --- versus $18\,000$ needed.
$\delta_{\mathrm{eff}}=\Theta(X^{-1/2})$ is an existence statement,
with the measured enumeration wall (work per certified in-window offset
scaling as $\mathrm{height}^{\theta}$, $\theta$ measured
$0.50$--$0.74$ against the theoretical $1/2$ [V]; windows at height
$10^{10}$ already contain zero enumerable hits at width $10^4$) as its
algorithmic shadow --- $\ge10^{100}$ operations per offset extrapolated
to $B=10^{200}$.

\emph{The dense frontier, $v\ge d$} [P at draft rigor]. When the value
set is window-dense, two elementary propositions close access anyway.
\emph{Isolation}: around a solution $u_0$ with $G(u_0)$ in the window,
any lattice displacement $\delta$ keeping $G(u_0+\delta)$ within a
width-$2Q$ window must satisfy $|\nabla G\cdot\delta|\lesssim Q$ and
$|\delta^{\mathsf T}H\delta|\lesssim Q$; the solution slab has
principal extents $Q\,X^{-(d-1)/d}$ and $Q^{1/2}X^{-(d-2)/(2d)}$ ---
at $d=4$, $X=10^{200}$: $10^{-145}$ and $10^{-47}$. Every extent is
below $\tfrac12$, so the only lattice point is $\delta=0$:
\emph{window solutions are pairwise isolated; there is no local
navigation between certified offsets --- dense existence is
non-constructive by geometry, not by hardness.} \emph{Orbit
accounting}: an infinite group $\Gamma$ of structured moves either
preserves $G$ (one offset per orbit), or preserves a factor --- so every
reachable value carries the planted divisor $a_0=A(u_0)$: enumerable
$a_0$ is a congruence certificate (covering economics), balanced $a_0$
is back in the isolation sliver --- or mixes generators of both
automorphism groups, whose coarse log-lattice of reachable factor pairs
holds $\sim(\log X)^2\approx2\cdot10^5$ points against a target of
relative width $10^{-195}$: expected hits $10^{-190}$. Any remaining
access route must solve $G(u)=m$ for \emph{prescribed} $m$ --- but a
product-form representation of $m$ \emph{is} a nontrivial factorization
of $m$, so a prescribed-value solver factors the scheme's own rough
complements: the extraction horn, not a new hypothesis. Disjoint-variable
products reduce directly to covering (the free small factor is the
certificate; \textsf{NormForm} reduces to \textsf{MultiPoly} through
exactly this route).

\paragraph{ExpFamily: closed by supply, not price [P, quantitative].}
Order conditions are the cheapest certificates in the grammar: $s$
divides $j\cdot a^n+c-q$ along an arithmetic progression of $n$ when
the discrete logarithm of $(q-c)j^{-1}$ lands in
$\langle a\rangle\bmod s$, at design cost $\ln\mathrm{ord}_s(a)$ ---
as low as $\ln3=1.10$ nats. But $s\mid a^{\mathrm{ord}}-1$ forces
$a^{\mathrm{ord}}>s$, so cheap certificates need small orders in large
$s$, and such $s$ supply only $\sim\pi(Q)\,\mathrm{ord}/s$ eligible
offsets each [e.g.\ $s=2801$, order 5: about 20]. Solving the resulting
purchase problem exactly --- sort all $s<10^5$ by price, buy cheapest
first until the pass bar of $1\,048$ offsets --- costs $1\,151$ nats
for every base $a\in\{2,3,5,7\}$, against the entire
$\ln B=460$-nat budget; the best \emph{budget-feasible} design
certifies $419$ offsets. The strongest non-covering family in the
grammar is $2.5\times$ over budget at the bar, ceiling $419$ against
$18\,000$.

\paragraph{Compose: empty at scale [P].} Two nontrivial family
constraints on the same $\N$ are contradictory: $\N=m^2-c_1=m'^2-c_2$
forces $(m-m')(m+m')=c_1-c_2$, so $m\le|c_1-c_2|$, while
$\N\sim10^{200}$ needs $m\sim10^{100}$ [V: ordinary shift pairs have
$O(1)$ solutions, largest $m=1073$ in the swept instances]. The only
surviving composition is family-plus-congruence --- a capped family
riding on ordinary covering, still under the bar. This explains
structurally the additivity the engine measured (no super-additive
interaction in any composed sample).

\paragraph{The economics, restated.} The generative tier
(\S\ref{sec:engine7}) forced a sharpening: with every mechanism charged
for \emph{all} offsets it claims and capped by the budget, a bulk
mechanism must deliver offsets at
$<\ln B/(3\sqrt Q)\approx0.44$ nats each --- and the honest comparison
is not covering's \emph{marginal} endgame ($1.01$ nats/offset) but its
\emph{average}: $437$ nats buy $10\,607$ certified offsets, $0.041$
nats each. Covering wins not at the margin but on average, because
$\ln r$ nats buy $\sim\pi(Q)/(r{-}1)$ offsets at once, whereas every
known algebraic mechanism buys $O(1)$ offsets per unit of entropy.

\section{Lemma 2: correlation exhaustion}\label{sec:lemma2}

\begin{proposition}[Shared divisors are small {[P]}]\label{prop:shared}
Let $\N$ be even, $q_1\ne q_2$ odd primes $<Q$, and $d>1$ with
$d\mid\N-q_1$ and $d\mid\N-q_2$. Then $d\mid q_1-q_2$, so every prime
factor of $d$ is $<Q$.
\end{proposition}
\begin{proof}
Subtract.
\end{proof}

Consequence: every divisor event coupling two complements is measurable
in the $\sigma$-algebra generated by small-prime divisibility --- which
is \emph{exactly} the information the cover fixes (classes) and the
sieve harvests ($k$-residues), and the engine consumes all of it: any
``$k$-side'' congruence trick selects a sub-progression the sieve
already identifies, and the survivor count is a sufficient statistic
(\S\ref{sec:engine}). Under H1, the residual lottery then factorizes,
$E_{\mathrm{eff}}$ decomposes as design cost plus $\sum_qp_q$, and that
is the covering objective [P conditional on H1]. Empirical support for
H1 at our scales is P5's $|\hat\Ee-\Ee|\le0.1$ over $1.2\cdot10^9$
candidates [V]; the corollary worth stating is that
\emph{luck-shaping} --- choosing $\N$ to make residual events
favourably dependent --- has no purchase beyond $0.1$ nats, because
dependence requires shared small divisors, which is covering by
definition.

\section{Lemma 3: the detector floor, posed as an open problem}\label{sec:lemma3}

The last escape is orthogonal to certificates: make \emph{testing}
cheaper than a modular exponentiation.

\begin{openproblem}[H2: sub-modexp compositeness detection on rough inputs]
\label{op:H2}
Let $n$ be uniformly random among integers in $[X,2X]$ with no prime
factor below $(\log X)^{O(1)}$, and let $M(b)$ be the bit cost of
$b$-bit multiplication. Does a randomized algorithm exist that decides
compositeness of $n$ with soundness $\ge2/3$ (one-sided; ``don't
know'' allowed) in time $o(M(\log X)\log X)$ --- asymptotically cheaper
than one modular exponentiation? Variants: (i) constant recall $r$ at
cost $c\cdot$modexp with $c/r<1$; (ii) amortized: certify a constant
fraction of the composites among $n_1,\dots,n_k$ in $o(k)$ modexps
total, beyond the small-divisor mass reachable by product
trees~\cite{Bernstein}; (iii) with polynomial advice depending on $X$
only.
\end{openproblem}

Known art sits firmly at the floor: trial division and batch gcd have
recall zero on rough inputs by construction; Jacobi symbols
equidistribute (soundness $o(1)$); Fermat, Miller--Rabin, and Frobenius
tests all cost $\Theta(M\log X)$; and no lower bound is known in any
realistic model --- even a conditional one in a restricted arithmetic
model would be new. The problem has a live industrial consumer: the
record engine spends $\ge95\%$ of its time on exactly this primitive
($\sim$17 tests per candidate, $e^{20.8}$--$e^{37.7}$ candidates per
campaign), as do the large distributed prime hunts; a detector at cost
$c$ and recall $r$ multiplies throughput by $1/(c+(1-r))$, and record
difficulty is $\log(\text{compute})$. A positive answer breaks the cost
floor of a whole genre of searches --- which is the more interesting
outcome, and why the problem is posed rather than assumed.

\section{The falsification engine}\label{sec:engine7}

A theorem whose proof is a case sweep earns its keep by inviting
attack. The engine's contract: a schema instance \emph{passes} ---
i.e.\ threatens the theorem --- iff simultaneously (1) it certifies
$\ge3\sqrt Q$ distinct prime offsets below $Q$ (clears the univariate
ceiling), with factor nontriviality spot-checked; (2) its cost per
offset, charged for \emph{every} offset claimed and budget-capped,
beats the covering benchmark; (3) its predicted effect would be
invisible in the banked density data (P5-consistency). Any pass is
audited by hand before it counts. The scorer re-derives every number
from primitives --- deterministic sieve, the pinned residual set,
symbolic factorization --- and \emph{proposals are data, never code};
fitness is capped, monotone in verified counts only, with no reward for
prose, novelty, or runtime.

\paragraph{Tier 0 (calibration).} The engine reproduces every P1--P4
hand number exactly, and surfaced two refinements worth keeping: the
covering endgame is $\approx1.0$ nats/offset at $Q=122\,000$ (an
earlier $2.5$-nat figure was context-bound), and the binding constraint
crosses over from budget to equilibrium between $Q=122\,000$ and
$Q=200\,000$ --- both now part of \S\ref{sec:regimes}.

\paragraph{The exhaustive tier.} $13\,661$ schema instances across the
full grammar: \textbf{zero passes}. Verdicts: $400$
congruence-composite instances scored \textsc{baseline} (inside the
paradigm --- the paradigm cannot escape itself); $12\,579$
\textsf{PolyImage} rejected on the value-set ceiling (best instance
remains $k{=}2$, $c{=}398$, $113$ offsets); $36$ \textsf{ExpFamily},
$324+4$ \textsf{MultiPoly}, $60$ \textsf{NormForm}, $198$
\textsf{Cyclotomic}, $60$ \textsf{Compose} rejected per the lemmas,
with all composition overlaps within $3\sigma$ Poisson of additive.
One methodological note kept for honesty: the first run produced $16$
spurious passes --- congruence-containing schemas leaking through the
escape filter with fresh-universe kill counts. The fix was semantic
(baseline family detection), not numeric; zero real passes before and
after. Near-miss board: \textsf{ExpFamily} at $3.1$--$4.2$ nats per
residual kill --- three times the covering endgame, exactly as P3
predicted --- then nothing below $10$.

\paragraph{The generative tier, and the referee it broke.} The
adversarial arm ran on a hosted evolutionary search service
(AlphaEvolve), under a threat model written in advance: evolved
candidates are \emph{programs}, so their only contract is to emit
schema JSON; they run sandboxed, never compute kills, never report
fitness; grading happens in a fresh interpreter. Within about ten
candidates the search drove apparent fitness $0.27\to0.46$ ---
not by finding mathematics, but by exploiting a hole in our scorer:
per-offset costs were billed only against the $867$ residuals of the
benchmark cover, so an \textsf{ExpFamily} design could advertise
$10\,416$ offsets costing $22\,900$ nats against a $460$-nat budget
and be credited anyway --- a $50\times$ overdraft. The audit protocol
caught it before any claim was made (a pass requires hand audit by
construction); the fix --- charge every claimed offset, cap coverage at
$\ln B/\mathrm{cost}$ --- re-scored the champion to $0.046$ and
\emph{strengthened} the theorem, yielding the budget constraint, the
$0.041$-nat average bar, and the purchase argument that closes
\textsf{ExpFamily} quantitatively (\S\ref{sec:lemma1}). The run then
converged: $41$ candidates, four families, zero passes, best corrected
fitness $0.0919$ --- the same near-miss family the exhaustive tier had
already ranked first, tuned to the edge of what order conditions allow,
stalled $53\times$ short of the bar.

\begin{table}[h]
\centering\small
\caption{Generative-tier empirical ceilings per family under corrected,
budget-aware accounting. ``eff.\ offsets'' is the budget-capped
certified count; the pass bar is $>1\,048$ offsets at $<0.44$
nats/offset, against covering's average $0.041$.}\label{tab:evolved}
\begin{tabular}{lrrrr}
\toprule
family & best fitness & eff.\ offsets & nats/offset & vs the bar\\
\midrule
\textsf{ExpFamily} & 0.092 & 209 & 1.10--2.20 & $53\times$ too expensive\\
\textsf{Compose} (2-deep) & 0.076 & 695 & 8.86 & $215\times$\\
\textsf{Cyclotomic} & 0.039 & 348 & 8.71 & $211\times$\\
\textsf{PolyImage} & 0.001 & 56 & 46.1 & $1118\times$\\
\bottomrule
\end{tabular}
\end{table}

Neither structural fact in Table~\ref{tab:evolved} is a search
artifact: no family clears both filter halves at once (coverage is
bought only at $8$--$9$ nats per offset; cheapness only at $O(1)$
offsets per condition), and the ordering matches Lemma~1's case
analysis exactly. The methodological finding generalises beyond this
thesis: \emph{adversarial search against a formal claim is most
valuable as a referee-hardening device, and should be reported that
way, not as a failed refutation} --- the designed behaviour of an audit
protocol in which candidates emit data, graders run fresh, and no pass
counts without human eyes.

\section{Door 1: the rounding gap}\label{sec:door1}

With mechanisms closed, the covering objective itself is the last
mathematical lever: how far from optimal are the covers? The answer
reframes the whole game.

\paragraph{The gap is everything.} The set-cover LP relaxation of the
benchmark instance ($Q=122\,000$; moduli fixed at the equilibrium set,
the $88$ odd primes $3..457$; one residue class to choose per modulus)
reports \emph{fractional coverage $100.00\%$} --- the classes carry
$\sum_r\pi(Q)/(r{-}1)\approx1.86\,\pi(Q)$ of kill mass, and fractional
routing covers every prime, at both $Q=122\,000$ and $Q=200\,000$. The
LP certifies nothing about integral assignments; what it proves is a
framing: \textbf{the entire game exponent is an integrality/rounding
gap}. Naive independent rounding leaves $e^{-1.86}\approx15.6\%$ of
primes uncovered; adaptive greedy leaves $7.56\%$; the fractional
optimum leaves $0$. Every absolute percent of residual removed is
$\approx2.7$ nats $\approx15\times$ less compute.

\paragraph{The landscape, measured.} Over the full $e^{436}$-point
assignment ensemble, random assignments have $|U|\sim(\mu=1\,377,
\sigma=21.0)$; the measured $\sigma$ sits \emph{below} the
independence prediction ($34.8$) --- negative correlations thin the
tail. A first-moment (annealed) threshold lands at $u^*\approx758$;
greedy-plus-descent achieves $867$, already $z=24.3$ of the $29.5$
available sigmas into the tail. Diagnostics: $19$ of $88$ moduli are
choice-tight; random-start coordinate ascent lands $90$--$110$ primes
\emph{worse} than greedy (adaptive construction is worth
$\approx2.4$ nats over naive local search); and pairwise class
agreement between independent solutions is $2.4\%$ --- chance level, a
fully degenerate landscape with no backbone for crossover methods to
graft.

\paragraph{Exactness, thrown at it.} The greedy point survived
everything: $41$ exact CP-SAT~\cite{ORtools} windows of $7$--$15$
moduli (warm-started, incumbent-floored) --- never beaten, always
matched, with $21$ of $22$ computed LP window bounds \emph{certifying}
window-optimality (the one uncertified window contains the smallest
moduli, the global LP's fractional looseness in miniature); exact
re-solves of the $30$ and then $44$ \emph{largest} moduli
simultaneously --- half the cover in one subproblem --- unchanged at
$867$; and belief-propagation-guided decimation (validated on a toy to
the exact optimum) landing at $874$, which its own local search cannot
improve --- a third mutually-isolated locally-exact basin
($867/874/968$), while the naive Bethe free energy slides toward the
fractional fantasy at low temperature, sharing the LP's blindness to
integrality (a survey-propagation treatment~\cite{MPZ} is the remaining
theory item). Running verdict: $867$ is at or very near the
\emph{quenched} optimum for this moduli set, the annealed $758$ likely
reflecting an annealed-versus-quenched (clustering) gap rather than
reachable headroom; the honest Door-1 prize is $0$--$2.6$ nats.

\paragraph{The open programme.} What has \emph{not} been tried is a
different ensemble: the covering constructions of Erd\H{o}s--Rankin
and FGKMT \cite{Rankin,FGKT,Maynard,FGKMT} are precisely a technology for
rounding fractional covers of primes with quantified loss --- layered
smooth-kill middle bands, second-moment control, a hypergraph-matching
endgame (cf.\ Moser--Tardos-style resampling~\cite{MT} for the
algorithmic version). Their loss is $o(1)$ in an asymptotic regime that
is not ours (their moduli reach $y/2$; ours cap at $457$), which is
exactly why this is a thesis-grade open question rather than a lookup:
\emph{determine the achievable residual fraction at
$\pi(Q)\approx10^4$ targets under budget $\theta(457)$.} Closing even
half the greedy-to-fractional gap is $e^{10}\approx22\,000\times$
compute --- the difference between a volunteer network and a
GPU-month for the $200\,000$ rung.

\section{Three doors, and the verdict}\label{sec:verdict}

Every attack generated in this programme reduces to one of three
irreducible doors; the closure argument is that there is no fourth.
\emph{(i)} Divisor information is fully harvested: $k$-side congruence
tricks select sub-progressions the sieve already identifies.
\emph{(ii)} Correlation requires a shared divisor, hence is covering
(Proposition~\ref{prop:shared}), and P5 bounds everything residual at
$0.1$ nats. \emph{(iii)} A compositeness certificate is planted (then
Lemma~1's accounting applies) or extracted (then it factors); and
certificate-free detection already runs at its conjectured floor (H2).

\begin{description}
\item[Door 1 --- rounding mathematics.] Prize $0$--$2.6$ nats on
current evidence, with the FGKMT-transfer ensemble as the genuine
unknown; every nat multiplies every future rung. Cheap, and the
covering monopoly theorem is what makes investing in it safe.
\item[Door 2 --- throughput.] The search is progress-free, stateless
per ticket, with millisecond verification: the best-shaped volunteer
workload there is, and the model on which the neighbouring
record-sports (GIMPS, PrimeGrid --- themselves covering problems in the
Sierpi\'nski--Riesel line) built everything. At the historical peak of
volunteer computing ($\sim$$2\times10^5$ of our fleets), the
$Q=200\,000$ rung falls in about two days; at $1\%$ of it, months.
The capital-intensive twin is a Montgomery-modexp ASIC rack
($\sim$$e^6$ over the fleet). The quantum door is closed at any
projected constants: Grover's halved exponent meets a
$\sim$$10^{11}$-fold oracle slowdown, crossover $\Ee^*\approx41$--$55$,
at or above every target of interest.
\item[Door 3 --- the paradigm.] This chapter. Status:
\textbf{proved on the swept grammar} [P at draft rigor] --- every
family closed by a proof of its own kind (value-set ceilings, cluster
bound, isolation plus orbit accounting, purchase argument, empty
intersection) --- \textbf{conjectural at the boundary}: H1 (supported
at $0.1$-nat resolution by $1.2\cdot10^9$ candidates), H2 (posed,
Open Problem~\ref{op:H2}), the $o(1)$ terms and equidistribution
inputs, and mechanisms outside $\mathcal L$ altogether --- neither a
planted divisor identity nor per-instance discovery --- of which none
is known.
\end{description}

The programme consequence is clean: with the covering monopoly a
result, record progress is governed by Door 1 and Door 2 --- and the
honest forecast for each game is Table~\ref{tab:walls}, unchanged by
anything a mechanism was hoped to deliver.

\chapter{Conclusions, open problems, and the honest epitaph}
\label{ch:conclusions}

\begin{chapterabstract}
We collect what the thesis established --- proved, measured, and
conjectured, marked apart --- restate the economics in one paragraph,
rank the open problems the work leaves behind, and record the
methodological practices that kept a record-driven, heavily automated
research programme honest.
\end{chapterabstract}

\section{What was established}\label{sec:established}

\paragraph{Proved.} $\g$ is unbounded, with explicit partitions, by
primorial obstruction plus Dirichlet (Theorem~\ref{thm:dirichlet});
any prime-offset cover is realisable by an actual Goldbach number
(Theorem~\ref{thm:coverdirichlet}); the maximal order of $\g$ is
$\gg\log\N\,\log_2\N\,\log_4\N/\log_3\N$ along explicit-partition
integers (Theorem~\ref{thm:lower}); shared divisors of two complements
are supported below $Q$ (Proposition~\ref{prop:shared}); and, at draft
rigor on the swept grammar, the Divisor-Paradigm Closure Theorem
(Theorem~\ref{thm:closure}): value-set ceilings, the cluster bound
$k_{\max}$, isolation and orbit accounting, the ExpFamily purchase
bound, and Compose emptiness --- covering is the unique economical
planted mechanism in $\mathcal L$.

\paragraph{Certified.} The record board of Table~\ref{tab:board}, and
behind it eleven new constructions (Table~\ref{tab:catalog} and the two
height records) each carrying
unconditional exclusion witnesses, a certified partition, dual-stack
verification, and manifests: the height records $\g=1\,157\,341$
($2\,480$ digits) and $\g=1\,134\,871$ ($2\,692$ digits, with its
$8\,970$ local gap certified, ratio $126.5$); the threshold descent to
$\N_{150}$ with $\g=104\,527$; the budget record $\g=119\,419$ under
$10^{199}$.

\paragraph{Measured.} The cost model $\Ee=|U|\cdot\mathrm{boost}/\ln\N$
to $\pm0.1$ nats over $1.2\cdot10^9$ candidates at five scales; the
frontier curves $\Ee(D)$, $\Ee(Q)$, and the million-desert wall
($10^{85}$ tests at $199$ digits against $\sim10^{6.5}$ existing
targets); the engine optimizations ($1.8\times$, $3\times$,
$430\times$); the compression ceilings (entropy-exact $k$-list
cancellation; $H^{0.35}$ conditioning saturation; the Johnson radius as
a practical wall at $1.028\,A_J$); the assignment landscape (fractional
coverage $100\%$; greedy at $z=24.3$; $867$ window-certified against
exact solvers up to half the cover); and the adversarial-search ceiling
per mechanism family (best corrected fitness $0.0919$, $53\times$ short).

\paragraph{Conjectured, named.} H1 (Hardy--Littlewood independence,
supported at $0.1$-nat resolution); H2
(Open Problem~\ref{op:H2}); the $o(1)$ terms and equidistribution
inputs of Lemma~1; out-of-grammar mechanisms; and the calibrated truths
--- $T(10^5)\approx10^{53}$, $R(10^{200})\approx1.8$ million --- that
give the records their scale bar.

\section{The three currencies, one paragraph}\label{sec:currencies}

A large least Goldbach summand is bought with three currencies ---
modulus digits, search throughput, and certification compute --- trading
at measured rates: about $2.5\times$ search per $7$ digits near
$190$, a digit per percent of cover quality, exponent
$e^{|U|\cdot\mathrm{boost}/\ln\N}$ for everything. Complete covers were
the gold standard and are obsolete; partial covers plus brute
progression scanning, checkpointed in git and certified twice by
independent code, are the entire present frontier. The truth for
$T(10^5)$ sits a factor $\sim$3 below the record in logarithm, and the
combined barrier analysis says only $10$--$20\%$ of that log-gap will
fall to compute and cover quality, with everything below
$\sim$$120$ digits protected by two independent entropy walls and a
closure theorem. The records are what they are: \emph{upper bounds,
manufactured, falling toward a floor they will never touch} --- a
sentence that belongs in every future README of this genre.

\section{Open problems, ranked}\label{sec:open}

Ranked by expected impact per unit effort, with the chapter that
motivates each.

\begin{enumerate}
\item \textbf{The rounding-gap programme}
(\S\ref{sec:door1}). Determine the achievable residual fraction for
one-class-per-modulus covers of $\pi(Q)\approx10^4$ primes under
budget $\theta(457)$ --- importing the Erd\H{o}s--Rankin/FGKMT layered
covering constructively at finite scale. Every absolute percent is
$\approx15\times$ compute; half the greedy-to-fractional gap is
$\approx22\,000\times$. Includes: certify the quenched threshold
(1RSB/survey propagation with complexity bookkeeping); lift-based
bounds (Sherali--Adams on the small-moduli block) where the plain LP is
provably vacuous; and the banded-ensemble construction with an exact
matching endgame.
\item \textbf{H2: sub-modexp compositeness detection}
(Open Problem~\ref{op:H2}). Either direction is a result: a detector
re-prices every prime-hunting search on earth; a lower bound in any
realistic model finishes Lemma~3.
\item \textbf{The two-sided construction} (\S\ref{sec:twosided}).
Intersect cover congruences with a certifiable-by-construction
complement family (Proth/Pocklington). Uncaps the height game ---
the only game with no mathematical wall --- from ECPP practice to
Proth-search scale. Key unknown: Proth-prime density inside
CRT-constrained progressions.
\item \textbf{CRT list-recovery beyond the Johnson radius}
(\S\ref{sec:subsetsum}--\ref{sec:toycrt}), for structured,
solution-abundant instances. The measured wall at $1.028\,A_J$ says
practice needs \emph{new decoding}, not better reduction.
\item \textbf{Modular subset-sum with choices at density
$\approx1$} when $2^{82}$ solutions exist (\S\ref{sec:subsetsum}).
Progress on either 4 or 5 collapses the sub-$100$ threshold verdict.
\item \textbf{Aggregate compositeness certificates.} Verification is
linear in the witness count; is there a sublinear batch certificate
that all $k$ complements are composite (resultant/product-tree/gcd
object verifiable faster than $k$ strong tests)? Even $5\times$
changes height-game economics; an impossibility result would be a
certificate-complexity theorem. Untouched.
\item \textbf{Matching lower bounds.} There is no theorem of the form
``every even $\N<B$ has $\g(\N)\le f(B)$'' beyond trivialities, even
conditionally. A GRH-conditional explicit bound, or a large-sieve bound
on how many residue classes a small $\N$ can occupy, would turn
Chapter~\ref{ch:limits}'s walls from forecast into mathematics.
\item \textbf{Complexity placement.} Is deciding
$\exists\,\text{even }\N<B:\g(\N)>Q$ (primality oracle given) NP-hard?
The reductions of Chapter~\ref{ch:limits} are evidence about
\emph{attacks}, not the problem; either answer locates the game on the
map.
\item \textbf{Neighbouring games, all virgin.} $S(q)$ exactly (the
smallest desert with prescribed least summand --- posed in 1989, no
constructive records exist); double deserts (force the two smallest
summands large); Lemoine deserts (odd $\N=p+2q$, force $\min p$);
ternary Goldbach deserts. The pipeline sets first-ever certified marks
in each essentially for free.
\item \textbf{Science from the exhaust.} The campaigns banked
$\sim$$10^9$ PRP tests of structured sparse sets at heights
$10^{150}$--$10^{2700}$ with full logs --- the largest such sample ever
produced. Testing Hardy--Littlewood $k$-tuple corrections against
fail-position statistics is a free empirical paper and a recalibration
of the $p_1$ model everything rests on.
\item \textbf{Formal verification.} The standalone verifier is
$\sim$200 lines of elementary arithmetic; a Lean/Coq proof of
``exit 0 $\Rightarrow$ $\g(\N)=q$'' would make these the first
machine-\emph{proof}-carrying records in the genre and close the last
trust gap.
\item \textbf{Correlation-aware covers.} Deliberately designing the
residual set for inadmissibility (minimal singular series) provably
raises the all-composite probability at fixed $|U|$; P5 bounds the
\emph{accidental} version at $0.1$ nats, so expect a modest constant
--- but nobody has optimised for it.
\end{enumerate}

For the engineering track, the costed roadmap (``Operation
Staircase'') stands ready: rungs at $Q=140$k, $150$k, $165$k, $175$k,
$200$k under $10^{200}$, each banked and dual-certified, with the final
$e^{5.5}$ wall closed by any one of volunteer computing, an ASIC rack,
or two nats of Problem 1 --- and a committed variant-ledger so parallel
workers scan disjoint space. The $2.5$-million height rung awaits only
its (priced, one-time) $3\,300$-digit certification.

\section{Methodological legacy}\label{sec:method}

Five practices earned permanent status, and are offered to any
record-driven computational programme:

\begin{enumerate}
\item \textbf{Verification-first.} Witnesses re-derived from published
data by standalone code, dual implementations for every certificate,
counts asserted against independent totals, manifests over everything.
The two false-PASS incidents (\S\ref{sec:certification}) were caught
by redundancy, not luck; no five-leg record was ever overturned.
\item \textbf{Adversarial audits, of ourselves first.} Planted
corruptions against the verifier; a hostile review of a contributed
blueprint before adoption; and the falsification engine's referee
broken by its own adversary within ten candidates --- the fix
strengthening the theorem it guarded. Never let a candidate report its
own fitness; never trust a pass without hands.
\item \textbf{Honest negatives, with statistics.} Hitless campaigns
reported with survival probabilities and per-slice model agreement;
one retraction (\S\ref{sec:subsetsum}) documented with its root cause
--- two bounds optimised at different parameters --- as a lesson, not
buried.
\item \textbf{Compute that survives its substrate.} Slices as
transactions, checkpoints in version control committed before results,
harvest-first workers: $\sim$25 container losses across the programme,
zero coverage lost.
\item \textbf{Models before campaigns.} No search was launched without
a validated density forecast, a $3\times$ budget, and written stopping
rules; every hit and every miss then updated the calibration. The
model never needed post-hoc adjustment --- which is the sentence that
makes all the others cheap.
\end{enumerate}

\section{Coda}\label{sec:coda}

The moral is a trade-off, not a paradox --- and by now it is a measured
one. Nature keeps the least Goldbach summand microscopically small, and
believes (\ref{eq:GLR}) it never grows much faster than
$(\log\N)^2\log\log\N$. Against that backdrop, deserts can be
manufactured: with fully static certificates at a thousand digits, or
with controlled search at a hundred and fifty, or past every certified
prime gap at twenty-five hundred --- each purchase itemised in a
currency whose exchange rates this thesis measured, and whose
counterfeits --- shortcuts through lattices, conditioning, or algebra
--- it priced honestly at infinity. What remains between the records
and the truth is now a short, explicit list of mathematical problems
rather than a fog. That list is the thesis's real deliverable: cute
gaps, and precisely posed ways to make them cuter.


\appendix
\chapter{Artifact map and verification runbook}
\label{app:artifacts}

Everything in this thesis is re-checkable from the artifact repository
(the \path{goldbach/} and \path{slop/goldbach/} trees). This appendix
maps claims to artifacts and gives the one-command verification for
each record. Environment as tested: Python~3.11 with \texttt{gmpy2},
\texttt{numpy}, \texttt{sympy} (and \texttt{python-flint} for the
lattice study), PARI/GP~2.15.4~\cite{PARI}; the \path{slop} verifier
runs on the Python standard library alone.

\section{Records and their verifiers}

Each record directory contains \path{record.json} (construction data:
cover pairs, $\N_0$, $M$, $k$, the decimal $\N$), the evidence file
(per-offset witnesses), the ECPP certificate(s) (\path{*.gp}), the
validation log, and a SHA-256 manifest.

In the table, G-\path{records/} abbreviates \path{goldbach/records/}
and S-\path{records/} abbreviates \path{slop/goldbach/records/}; the
GPU/goldbach-pipeline records verify with
\path{goldbach/verify_record.py} (the mega-desert with
\path{goldbach/verify_megagap.py}), the CPU/slop-pipeline records with
the stdlib-only \path{slop/goldbach/records/verify_record.py}.

\begin{center}\footnotesize
\begin{tabular}{lll}
\toprule
record & directory & stack\\
\midrule
$\g=104\,527$, 150d (T) & G-\path{records/threshold_150digit_g104527} & G\\
$\g=100\,297$, 195d & G-\path{records/threshold_195digit_g100297} & G\\
$\g=107\,719$, 197d & G-\path{records/dual_197digit_g107719} & G\\
$\g=112\,249$, 200d & G-\path{records/budget_1e200_200digit_g112249} & G\\
$\g=1\,134\,871$, 2692d & G-\path{records/megagap_2692digit_g1134871} & G (megagap)\\
$\g=1\,157\,341$, 2480d (H) & S-\path{records/g1} & S\\
$\g=110\,917$, 199d & S-\path{records/g23} & S\\
$\g=102\,337$, 193d & S-\path{records/t100k} (round 1) & S\\
$\g=109\,357$, 186d & S-\path{records/t100k} (round 2) & S\\
$\g=101\,149$, 179d & S-\path{records/t100k} (round 3) & S\\
$\g=119\,419$, 199d (R) & S-\path{records/r200} & S\\
\bottomrule
\end{tabular}
\end{center}

Typical invocations:
\begin{verbatim}
python3 goldbach/verify_record.py \
    goldbach/records/threshold_150digit_g104527/record.json /tmp/ev.json
python3 goldbach/verify_megagap.py       # full megagap re-check + manifest
python3 slop/goldbach/records/verify_record.py \
    slop/goldbach/records/g1/record_g1.json
\end{verbatim}
The verifiers reconstruct $\N$ by CRT, re-derive every witness below
$\g$ (asserting the scanned count against an independent $\pi(\g)$),
re-verify the ECPP chain from scratch including the binding of its top
integer to $\N-\g$, and check the record inequalities; exit code $0$
means all legs pass. ECPP certificates additionally validate under
PARI's \texttt{primecertisvalid}, and the independent projective
checker is \path{goldbach/ecpp_check.py}
(\path{slop/.../g23/ecpp_verify.py} for the CPU-session records, which
also carry APR-CL logs and \path{crosscheck_*.gp} full-scan scripts).

\section{Search, cover, and analysis tools}

Paths below are relative to \path{goldbach/}.

\begin{center}\footnotesize
\begin{tabular}{p{0.40\textwidth}p{0.50\textwidth}}
\toprule
tool & role (chapter)\\
\midrule
\path{cover.py}, \path{anneal.py} & greedy weighted covers; SA refinement (\ref{ch:records})\\
\path{gen_variants.py} & residue-swap variant catalogs (\ref{ch:records})\\
\path{search.py} & sieved progression scan; \texttt{--sort-tests}, \texttt{--skip-frac} (\ref{ch:records})\\
\path{gpu/engine150.py}, \path{gpu/fleet.py} & CUDA engine and fleet driver (\ref{ch:records})\\
\path{verify_record.py}, \path{package_record.py} & witness generation; record packaging (\ref{ch:records})\\
\path{joint.py} & joint representative/coverage optimizer, $k$-list meet-in-the-middle (\ref{ch:limits})\\
\path{exp1m.py} & million-desert frontier and wall (\ref{ch:costmodel})\\
\path{entropy_frontier.py}, \path{realizable_frontier.py}, \path{frontier_table.py} & set-valued conditioning study (\ref{ch:limits})\\
\path{verify_barrier.py} & existence-vs-decoding coupling table (\ref{ch:limits})\\
\path{phd/toy_crt/decoder.py} & LLL vs Johnson radius study (\ref{ch:limits})\\
\path{phd/wp8/lp_bound.py}, \path{phd/wp8/pack_probe.py} & LP vacuousness; packing probe (\ref{ch:closure})\\
\path{phd/wp9/engine/} & closure-theorem engine: \path{tier0.py}, \path{exhaustive.py}, \path{scorer.py}, harness docs (\ref{ch:closure})\\
\path{phd/wp9/door1/} & rounding-gap instance, CP-SAT LNS, BP/SP, verdicts (\ref{ch:closure})\\
\bottomrule
\end{tabular}
\end{center}

Constructions referenced in Chapter~\ref{ch:limits} are committed
under \path{goldbach/constructions/}: \path{m211_x196.json},
\path{m219_x197.json}, \path{m221_x197.json},
\path{wall_1M_199d.json}, and \path{blueprint_1M_1250d.json}, each
re-verified end to end by its generator.

\section{Research notes underlying each chapter}

The narrative documents merged from the research sessions, kept as
primary sources (paths again relative to \path{goldbach/}):
\path{RESULTS.md} (record announcements and statistics),
\path{SOK_GOLDBACH_GAPS.md} (systematization, walls),
\path{FRONTIER.md} (set-valued conditioning and sub-100 wall),
\path{IDEAS.md} and \path{IDEAS100.md} (the idea ledgers with
kill-status), \path{PROPOSAL_NEXT_CAMPAIGNS.md} (the staircase),
\path{phd/RESEARCH_LOG.md} (the dated log of campaigns, audits, and
merges), the memo \path{phd/wp3_small_representative_memo.md}
(subset-sum barrier, including the v1 retraction),
\path{phd/wp7_probes.md} (the five probes),
\path{phd/wp8_beyond_covering.md} (doors; rung ladder; LP result),
\path{phd/wp9/wp9_closure_theorem.md} (the theorem memo, \S1--14),
\path{phd/wp9/h2_open_problem.md} (H2, posting-ready), and
\path{phd/wp9/door1/} (landscape diagnostics and gate verdicts).
Campaign drivers and their committed checkpoint logs live under
\path{phd/campaign*/} and, for the CPU session,
\path{slop/goldbach/records/work/}.

\section{Independence disclosure}

Two verification stacks were developed in separate sessions with no
shared code: the \path{goldbach} stack (\path{verify_record.py} +
\path{ecpp_check.py} + PARI) and the \path{slop} stack
(stdlib-only \path{verify_record.py} + \path{ecpp_verify.py} + APR-CL
+ PARI full scan). Every headline record has passed both. The
adversarial audit transcript (planted corruptions, all detected) and
the two false-PASS incidents with their root causes are recorded in
\path{goldbach/SOK_GOLDBACH_GAPS.md} \S4 and
\path{goldbach/phd/RESEARCH_LOG.md}.


% ---------------------------------------------------------------- bibliography
\begin{thebibliography}{99}

\bibitem{OeS} T.~Oliveira e Silva, S.~Herzog, S.~Pardi, \emph{Empirical
verification of the even Goldbach conjecture and computation of prime
gaps up to $4\cdot10^{18}$}, Math.\ Comp.\ \textbf{83} (2014), 2033--2060.

\bibitem{Richstein} J.~Richstein, \emph{Verifying the Goldbach conjecture
up to $4\cdot10^{14}$}, Math.\ Comp.\ \textbf{70} (2001), 1745--1749.

\bibitem{HL} G.~H.~Hardy, J.~E.~Littlewood, \emph{Some problems of
`Partitio Numerorum' III}, Acta Math.\ \textbf{44} (1923), 1--70.

\bibitem{GLvdtR} A.~Granville, J.~van de Lune, H.~J.~J.~te Riele,
\emph{Checking the Goldbach conjecture on a vector computer}, in: Number
Theory and Applications (R.~A.~Mollin, ed.), NATO ASI Ser.~C
\textbf{265}, Kluwer, Dordrecht, 1989, 423--433.

\bibitem{OEIS} N.~J.~A.~Sloane (ed.), \emph{The On-Line Encyclopedia of
Integer Sequences}, sequences A020481 (least Goldbach summand) and
A025018/A025019 (records), \url{https://oeis.org}.

\bibitem{JB} Zs.~Juh\'asz, M.~Bartalos, \emph{Predicting the size ranking
of minimal primes in the generalised Goldbach partitions},
arXiv:2510.21870 (2025).

\bibitem{Dirichlet} P.~G.~L.~Dirichlet, \emph{Beweis des Satzes, dass
jede unbegrenzte arithmetische Progression\dots\ unendlich viele
Primzahlen enth\"alt}, Abh.\ K.\ Preuss.\ Akad.\ Wiss.\ (1837), 45--81.

\bibitem{Linnik} U.~V.~Linnik, \emph{On the least prime in an arithmetic
progression I, II}, Mat.\ Sbornik \textbf{15(57)} (1944), 139--178,
347--368.

\bibitem{Rankin} R.~A.~Rankin, \emph{The difference between consecutive
prime numbers}, J.\ London Math.\ Soc.\ \textbf{13} (1938), 242--247.

\bibitem{Erdos1950} P.~Erd\H{o}s, \emph{On integers of the form $2^k+p$
and some related problems}, Summa Brasil.\ Math.\ \textbf{2} (1950),
113--123.

\bibitem{FGKT} K.~Ford, B.~Green, S.~Konyagin, T.~Tao, \emph{Large gaps
between consecutive prime numbers}, Ann.\ of Math.\ \textbf{183} (2016),
935--974.

\bibitem{Maynard} J.~Maynard, \emph{Large gaps between primes}, Ann.\ of
Math.\ \textbf{183} (2016), 915--933.

\bibitem{FGKMT} K.~Ford, B.~Green, S.~Konyagin, J.~Maynard, T.~Tao,
\emph{Long gaps between primes}, J.\ Amer.\ Math.\ Soc.\ \textbf{31}
(2018), 65--105.

\bibitem{GapList} The Prime Gap List project (successor of T.~R.~Nicely's
tables), \emph{Tables of first known occurrence prime gaps},
\url{https://primegap-list-project.github.io/}.

\bibitem{PSW} C.~Pomerance, J.~L.~Selfridge, S.~S.~Wagstaff~Jr., \emph{The
pseudoprimes to $25\cdot10^9$}, Math.\ Comp.\ \textbf{35} (1980),
1003--1026.

\bibitem{BW} R.~Baillie, S.~S.~Wagstaff~Jr., \emph{Lucas pseudoprimes},
Math.\ Comp.\ \textbf{35} (1980), 1391--1417.

\bibitem{GoldwasserKilian} S.~Goldwasser, J.~Kilian, \emph{Almost all
primes can be quickly certified}, Proc.\ 18th ACM Symposium on Theory of
Computing (1986), 316--329.

\bibitem{AtkinMorain} A.~O.~L.~Atkin, F.~Morain, \emph{Elliptic curves
and primality proving}, Math.\ Comp.\ \textbf{61} (1993), 29--68.

\bibitem{APRCL} L.~M.~Adleman, C.~Pomerance, R.~S.~Rumely, \emph{On
distinguishing prime numbers from composite numbers}, Ann.\ of Math.\
\textbf{117} (1983), 173--206; H.~Cohen, A.~K.~Lenstra, \emph{Primality
testing and Jacobi sums}, Math.\ Comp.\ \textbf{42} (1984), 297--330.

\bibitem{Pocklington} H.~C.~Pocklington, \emph{The determination of the
prime or composite nature of large numbers by Fermat's theorem}, Proc.\
Cambridge Philos.\ Soc.\ \textbf{18} (1914), 29--30.

\bibitem{Proth} F.~Proth, \emph{Th\'eor\`emes sur les nombres premiers},
C.\ R.\ Acad.\ Sci.\ Paris \textbf{87} (1878), 926.

\bibitem{PARI} The PARI Group, \emph{PARI/GP version 2.15.4},
\url{https://pari.math.u-bordeaux.fr}.

\bibitem{GSS} V.~Guruswami, A.~Sahai, M.~Sudan, \emph{``Soft-decision''
decoding of Chinese Remainder codes}, Proc.\ 41st IEEE Symposium on
Foundations of Computer Science (FOCS 2000), 159--168.

\bibitem{HG} N.~Howgrave-Graham, \emph{Approximate integer common
divisors}, Cryptography and Lattices (CaLC 2001), Lecture Notes in
Comput.\ Sci.\ \textbf{2146}, Springer, 51--66.

\bibitem{LO} J.~C.~Lagarias, A.~M.~Odlyzko, \emph{Solving low-density
subset sum problems}, J.\ Assoc.\ Comput.\ Mach.\ \textbf{32} (1985),
229--246.

\bibitem{CJLOSS} M.~J.~Coster, A.~Joux, B.~A.~LaMacchia, A.~M.~Odlyzko,
C.-P.~Schnorr, J.~Stern, \emph{Improved low-density subset sum
algorithms}, Comput.\ Complexity \textbf{2} (1992), 111--128.

\bibitem{Wagner} D.~Wagner, \emph{A generalized birthday problem},
Advances in Cryptology --- CRYPTO 2002, Lecture Notes in Comput.\ Sci.\
\textbf{2442}, Springer, 288--304.

\bibitem{Bernstein} D.~J.~Bernstein, \emph{Fast multiplication and its
applications}, in: Algorithmic Number Theory (J.~P.~Buhler,
P.~Stevenhagen, eds.), MSRI Publications \textbf{44}, Cambridge
University Press, 2008, 325--384.

\bibitem{Grover} L.~K.~Grover, \emph{A fast quantum mechanical algorithm
for database search}, Proc.\ 28th ACM Symposium on Theory of Computing
(1996), 212--219.

\bibitem{GE} C.~Gidney, M.~Eker\aa, \emph{How to factor 2048 bit RSA
integers in 8 hours using 20 million noisy qubits}, Quantum \textbf{5}
(2021), 433.

\bibitem{GranvilleCramer} A.~Granville, \emph{Harald Cram\'er and the
distribution of prime numbers}, Scand.\ Actuar.\ J.\ \textbf{1995:1},
12--28.

\bibitem{MPZ} M.~M\'ezard, G.~Parisi, R.~Zecchina, \emph{Analytic and
algorithmic solution of random satisfiability problems}, Science
\textbf{297} (2002), 812--815.

\bibitem{MT} R.~A.~Moser, G.~Tardos, \emph{A constructive proof of the
general Lov\'asz local lemma}, J.\ Assoc.\ Comput.\ Mach.\ \textbf{57}
(2010), Article 11.

\bibitem{ORtools} L.~Perron, F.~Didier, Google OR-Tools
CP-SAT solver, \url{https://developers.google.com/optimization}.

\end{thebibliography}

\end{document}
